\begin{textsample}{POS Dim 2 – 2021\_c6 – Score 3.00 – 2021\_c6/enade\_2021\_filo\_baa597e9cf.txt}  \label{ex:f2_pos_176}
CAPÍTULO 6 ANÁLISE TÉCNICA DA PROVA Este capítulo tem por objetivo apresentar o desempenho dos estudantes concluintes de Filosofia ( Licenciatura/Bacharelado ) no Enade/2021.
Para isso, foram calculadas as estatísticas básicas da prova em seu todo ( seção 6.1.1 ), bem como as estatísticas relacionadas aos Componentes de Formação Geral ( seção 6.1.2 ) e de Conhecimento Específico da Área ( seção 6.1.3 ).
Dadas as suas características, foram analisadas, em separado, as questões objetivas ( seção 6.2 ) e as questões discursivas ( seção 6.3 ).
Tomando -se como base as duas questões discursivas do Componente de Formação Geral, nas seções 6.3.1.6 a 6.3.1.8, são apresentados os resultados e os comentários sobre a correção das respostas em relação à Língua Portuguesa.
Para as questões objetivas, foram disponibilizados os índices de Facilidade e de Discriminação Ponto-Bisserial, também em separado, para os Componentes de Formação Geral ( seção 6.2.1 e 6.2.2 ) e de Conhecimento Específico ( seção 6.2.3 e 6.2.4 ).
De cada componente, uma das questões foi escolhida para exemplificar a análise gráfica, relacionando -se as alternativas escolhidas pelos estudantes ( inclusive o gabarito ) com o número de acertos no componente.
Nos Anexos IA e IB, apresenta -se a íntegra da análise gráfica para todas as questões objetivas, respectivamente à Licenciatura e ao Bacharelado.
Para cada uma das questões discursivas, os conteúdos dos tipos mais comuns de respostas dos estudantes são apresentados e comparados com o padrão de resposta esperado ( ver Anexos IXA e IXB com o padrão de respostas ).
Nas tabelas, constam as seguintes estatísticas das notas¹⁹ : média do desempenho na prova, erro padrão da média, desvio padrão, nota mínima, mediana e nota máxima para cada um de seus componentes.
Tais estatísticas contemplam o total de estudantes concluintes da Área de Filosofia ( Licenciatura/Bacharelado ) inscritos e presentes à prova do Enade/2021, tendo em vista agregações, ou por Grandes Regiões e o país como um todo, segundo Habilitação e Modalidade de Ensino, ou por Categoria Administrativa, Organização Acadêmica da IES.
Em relação aos gráficos de distribuição de notas, o intervalo considerado foi de 10 unidades, aberto à esquerda e fechado à direita, com exceção do primeiro intervalo, [ 0 ; 10 ], fechado em ambos os extremos.
Para os gráficos de distribuição das notas das questões discursivas, foram consideradas mais duas categorias : questão em branco²⁰ e nota zero. ¹⁹ Uma definição dessas estatísticas pode ser encontrada no Glossário. ²⁰ Nesse grupo estão incluídas também as respostas classificadas como nulas ou desconsideradas. 6.1 ESTATÍSTICAS BÁSICAS DA PROVA Esta seção apresenta estatísticas selecionadas e histogramas da nota geral ( 6.1.1 ) e de cada componente : Formação Geral ( 6.1.2 ) e Conhecimento Específico ( 6.1.3 ).
São também apresentadas estatísticas selecionadas de subpopulações, caracterizadas por Grande Região, Categoria Administrativa, Organização Acadêmica e Modalidade de Ensino. 6.1.1 ESTATÍSTICAS BÁSICAS GERAIS Na Tabela 6.1, são apresentadas as Estatísticas Básicas da prova, por Grande Região segundo Habilitação e Modalidade de Ensino, dos estudantes concluintes de Filosofia ( Licenciatura/Bacharelado ).
As análises serão feitas primeiramente para os cursos de Licenciatura e depois para os de Bacharelado.
Cada habilitação é desagregada internamente em cursos Presenciais e cursos a Distância.
A Média das notas da prova, como um todo, dos alunos de Licenciatura ( nas seções seguintes serão analisados os Componentes de Formação Geral e de Conhecimento Específico ), foi 48,8, sendo que os estudantes da região Norte obtiveram a Média mais baixa ( 42,4 ), e os da região Sudeste obtiveram a Média mais alta ( 51,0 ).
As demais Médias foram : 45,6, na região Nordeste, 49,5, na região Sul e 48,0, na região Centro-Oeste.
O Desvio padrão para o Brasil, como um todo, foi 16,5, sendo os maiores Desvio padrão encontrados na região Centro-Oeste ( 17,3 ), e o menor, na região Norte ( 15,2 ), indicando uma dispersão um pouco menor das notas dessa última região.
A região que obteve a maior nota Máxima foi a Sul ( 91,1 ), ao passo que a região que atingiu a menor nota Máxima foi a Norte ( 75,8 ).
A Mediana do Brasil, como um todo, foi 49,6, sendo a maior Mediana obtida na região Sudeste ( 52,2 ), e a menor obtida na região Norte ( 44,5 ).
A nota Mínima do Brasil ( zero ) foi obtida em todas regiões, com exceção da região Nordeste onde foi obtida a nota Mínima 7,4.
Considerando -se as notas, como um todo dos alunos de Licenciatura segundo Grande Região, observa -se que existe diferença estatisticamente significativa ao nível de 95 \% entre as Médias obtidas na região Norte e nas regiões Sudeste, Sul e Centro-Oeste, assim como entre a região Nordeste e as regiões Sudeste e Sul.
A Média das notas para os alunos de cursos Presenciais de Licenciatura, foi 51,4, sendo que os alunos da região Norte obtiveram a Média mais baixa ( 42,4 ), e os da região Sudeste obtiveram a Média mais alta ( 56,5 ).
As demais Médias foram : 48,2 na região Nordeste, 55,2 na região Sul e 47,1 na região Centro-Oeste.
O Desvio padrão das Licenciaturas Presenciais foi 16,1, sendo o menor obtido na região Nordeste 14,6 e o maior na região Centro-Oeste ( 17,7 ).
A maior Mediana foi a da região Sudeste ( 58,9 ) e a menor foi a da região Norte ( 44,5 ).
A maior nota Máxima foi 90,8, obtida na região Sudeste, registrada como a nota Máxima de todas as Licenciaturas.
As notas Mínimas foram diferentes de zero nas regiões Nordeste ( 9,3 ) e Sul ( 11,1 ).
Na Modalidade de Educação Presencial de Licenciatura, não existe diferenças, estatisticamente significativa ao nível de 95 \%, entre as Médias obtidas nas Grandes Regiões Norte e Centro-Oeste, Centro-Oeste e Nordeste, e Sudeste e Sul.
A Média das notas da prova dos alunos de cursos a Distância de Licenciatura, foi 45,9, mais baixa do que os alunos de cursos Presenciais ( 51,4 ).
Somente na região Norte não foram oferecidos cursos a Distância de Licenciatura.
A Média mais alta foi obtida na região Centro-Oeste ( 54,4 ) e a mais baixa foi obtida na região Nordeste ( 34,0 ).
O Desvio padrão das Licenciaturas a Distância foi 16,6, sendo o menor obtido na região Centro-Oeste ( 12,5 ) e o maior na região Sul ( 16,5 ).
A maior nota Máxima foi 91,1, obtida na região Sul, registrada como a nota Máxima de todas as Licenciaturas.
As notas Mínimas foram diferentes de zero nas regiões Nordeste ( 7,4 ) e Centro-Oeste ( 27,1 ).
Na Modalidade de Educação a Distância de Licenciatura, não existe diferenças, estatisticamente significativa ao nível de 95 \%, entre as Médias obtidas nas Grandes Regiões Sudeste e Sul.
Cumpre notar que ainda que as notas dos cursos Presenciais sejam comparáveis com as dos cursos a Distância dentro de uma mesma Habilitação, não é possível comparar as notas obtidas pelos alunos de Licenciatura e de Bacharelado, já que as provas são, em parte, diferentes.
A Média das notas da prova, como um todo dos alunos de Bacharelado ( nas seções seguintes serão analisados os Componentes de Formação Geral e de Conhecimento Específico ), foi 46,1, sendo que os estudantes da região Norte obtiveram a Média mais baixa ( 28,1 ), e os da região Sul obtiveram a Média mais alta ( 49,3 ), sendo maior do que a Média do Brasil.
As demais Médias foram : 40,4, na região Nordeste, 47,1, na região Sudeste e 49,2 na região Centro-Oeste ( onde também foi maior do que a Média do Brasil ).
O Desvio padrão para o Brasil, como um todo, foi 16,1, sendo o maior Desvio padrão encontrado nas regiões Sudeste e Centro-Oeste ( 16,1 ), e o menor, na região Norte ( 10,9 ), indicando uma dispersão um pouco menor das notas dessa última região.
A região que obteve a maior nota Máxima foi a Centro-Oeste ( 85,6 ), ao passo que a região que atingiu a menor nota Máxima foi a Norte ( 53,7 ).
A Mediana do Brasil, como um todo, foi 46,5, sendo a maior Mediana obtida na região Sul ( 50,0 ), e a menor obtida na região Norte ( 27,3 ).
A nota Mínima do Brasil foi obtida na região Sudeste, onde foi zero.
A maior nota Mínima, entre as cinco Grandes Regiões, foi 18,1, na região Centro-Oeste, as demais regiões as notas Mínimas foram : 1,7 na região Norte, 10,6 na região Nordeste e 11,3 na região Sul.
Considerando -se as notas de todo o Brasil dos alunos de Bacharelado, segundo Grande Região, observa -se que não existe diferença estatisticamente significativa ao nível de 95 \% entre as Médias, obtidas nas regiões Sudeste, Sul e Centro-Oeste.
Já as duas menores Médias, obtidas nas regiões Norte e Nordeste, têm diferença estatisticamente entre si e em relação às demais Médias.
A Média das notas da prova como um todo para os alunos de cursos Presenciais de Bacharelado, foi 46,1, sendo que os alunos da região Norte obtiveram a Média mais baixa ( 28,1 ), e os da região Sul obtiveram a Média mais alta ( 49,6 ).
As demais Médias foram : 40,4 na região Nordeste, 48,0 na região Sudeste e 48,4 na região Centro-Oeste.
A maior nota Máxima dos cursos presenciais de Bacharelado foi a da região Centro-Oeste ( 85,6 ) e a menor foi a da região Norte ( 53,7 ).
A maior Mediana obtida na região Sul ( 50,8 ), e a menor obtida na Norte ( 27,3 ).
As notas Mínimas nos cursos Presenciais segundo região, foram : Norte ( 1,7 ), Nordeste ( 10,6 ), Sudeste ( 0,0 ), Sul ( 11,3 ) e Centro-Oeste ( 18,1 ).
Considerando -se as notas segundo Grande Região, não existe diferença estatisticamente significativa ao nível de 95 \% entre as Médias das regiões Sudeste, Sul e Centro-Oeste.
Somente nas regiões Norte e Nordeste não foram oferecidos cursos a Distância de Bacharelado A Média mais alta foi obtida na região Centro-Oeste ( 51,2 ) e a mais baixa foi obtida na região Sudeste ( 40,6 ).
O Desvio padrão dos Bacharelados a Distância foi 15,4, sendo o menor obtido na região Sul ( 13,3 ) e o maior na região Centro-Oeste ( 16,7 ).
A maior nota Máxima foi 84,3, obtida na região Sul, registrada como a nota Máxima de todos os Bacharelados.
As notas Mínimas segundo região, foram : Sudeste ( 12,7 ), Sul ( 20,5 ) e Centro-Oeste ( 22,2 ).
Na Modalidade de Educação a Distância de Bacharelado, não existe diferenças, estatisticamente significativa ao nível de 95 \%, entre as Médias obtidas nas Grandes Regiões Sul e Centro-Oeste.
O comportamento das notas dos estudantes de todo o Brasil pode ser observado no Gráfico 6.1, que apresenta um histograma com a distribuição das mesmas, segundo a Habilitação.
Essa é uma distribuição unimodal para os alunos de cursos de Licenciatura e de Bacharelado, com ambas as modas no intervalo ( 50 ; 60 ].
Na Tabela 6.2, são apresentadas informações referentes à Média da nota final, desagregadas por Categoria Administrativa, por Organização Acadêmica, segundo a Habilitação.
A Média das notas da prova, como um todo, dos alunos de Licenciatura, foi 48,8 ( Tabela 6.1 ).
Em relação a Categoria Administrativa, os estudantes de Licenciatura das IES Públicas obtiveram Média mais alta ( 49,3 ) que a Média nacional, e os das IES Privadas obtiveram Média mais baixa ( 48,3 ).
Observa -se que não existe diferença estatisticamente significativa entre as Médias das notas das IES Públicas e as das Privadas.
A diferença entre as Médias das regiões Sul e Norte ( 8,6 ), a maior e a menor Média, é superior à diferença entre a das IES Públicas e Privadas ( 1,0 ), caracterizando uma maior diversidade regional do que administrativa.
O Desvio padrão para as IES Públicas ( 16,8 ) foi superior ao do Brasil, como um todo ( 16,5 ), indicando uma dispersão um pouco maior das notas nesta Categoria Administrativa.
No tocante à Organização Acadêmica, as Universidades ( 49,4 ) obtiveram Média mais alta que a nacional ( 48,8 ).
A Média dos Centros Universitários ( 47,4 ) e a das Faculdades ( 48,4 ) foram menores do que a nacional.
Constata -se que não existe diferença estatisticamente significativa, ao nível de 95 \%, entre as Médias dos Centros Universitários e Faculdades.
A Organização Acadêmica CEFET/IF não teve cursos de Licenciatura em Filosofia avaliados.
A Média das notas da prova, como um todo, dos alunos de Bacharelado, foi 46,1 ( Tabela 6.1 ).
Em relação a Categoria Administrativa, os estudantes de Bacharelado das IES Públicas obtiveram Média mais alta ( 51,4 ) que a Média nacional, e os das IES Privadas obtiveram Média mais baixa ( 42,5 ).
Observa -se que existe diferença estatisticamente significativa entre as Médias das notas das IES Públicas e as das Privadas.
A diferença entre as Médias das regiões Sul e Norte ( 21,2 ), a maior e a menor Média, é superior à diferença entre a das IES Públicas e Privadas ( 8,9 ), caracterizando -se uma menor diversidade administrativa do que regional.
O Desvio padrão para as IES Públicas ( 15,6 ) foi inferior ao do Brasil, como um todo ( 16,1 ), indicando uma dispersão um pouco menor das notas nesta Categoria Administrativa.
No tocante a Organização Acadêmica, as Universidades obtiveram Média mais alta que a nacional ( 48,3 ).
A Média dos Centros Universitários e a das Faculdades foram menores do que a nacional ( 43,5 e 39,9, respectivamente ).
Constata -se que não existe diferença estatisticamente significativa, ao nível de 95 \%, entre as Médias dos Centros Universitários e Faculdades.
A Organização Acadêmica CEFET/IF não teve cursos de Bacharelado em Filosofia avaliados. 6.1.2 ESTATÍSTICAS BÁSICAS NO COMPONENTE DE FORMAÇÃO GERAL Na Tabela 6.3, são apresentadas as Estatísticas Básicas em relação ao componente da prova que avalia a Formação Geral dos estudantes concluintes.
A nota de Formação Geral é obtida levando -se em conta os acertos nas oito questões objetivas para ambas habilitações, sendo a nota de cada questão discursiva obtida levando -se em conta 80 \% da nota relativa ao conteúdo e 20 \% da nota relativa ao desempenho linguístico.
Os estudantes de Licenciatura, de todo o Brasil, obtiveram desempenho médio de 38,6.
Quanto à variabilidade, o Desvio padrão das notas dos estudantes de Licenciatura do Brasil foi 15,2.
A maior Média foi obtida na região Sudeste ( 40,9 ) e a menor, na região Norte ( 33,7 ).
As demais Médias foram : 35,6 na região Nordeste, 38,5 na região Sul e 38,7 na região Centro-Oeste.
Já o maior Desvio padrão foi obtido na região Sudeste ( 15,5 ), e o menor, na região Nordeste ( 14,3 ).
Os demais desvios foram : 15,4, na região Norte, 15,0, na região Sul e 15,1, na região Centro-Oeste.
A maior nota no Componente de Formação Geral da prova do Enade/2021 foi 88,1, obtida por, pelo menos, um estudante na região Sudeste.
A menor nota Máxima foi obtida na região Norte ( 72,2 ).
A Mediana do Brasil, como um todo, foi 38,6, sendo a menor Mediana encontrada na região Norte ( 34,4 ), e a maior, na região Sudeste ( 41,2 ).
A nota Mínima nessa parte foi zero em todas as regiões.
Considerando -se as notas, segundo Grande Região, observa -se que não existe diferença estatisticamente significativa entre as Médias das notas do Componente de Formação Geral obtidas nas regiões Norte e Nordeste e entre Sudeste e Centro-Oeste.
As diferenças entre as demais Médias são estatisticamente significativas.
Já os alunos de cursos Presenciais de Licenciatura obtiveram desempenho médio de 39,6.
Quanto à variabilidade, o Desvio padrão das notas do conjunto dos estudantes das Licenciaturas presenciais foi 15,1.
A maior Média foi obtida na região Sudeste ( 43,2 ), e a menor, na região Norte ( 33,7 ).
As demais Médias foram : 37,1 na região Nordeste, 42,0 na região Sul e 37,5 na região Centro-Oeste.
O maior Desvio padrão foi obtido na região Norte ( 15,4 ), e o menor, nas regiões Nordeste ( 14,1 ).
Os demais desvios foram : 15,1 na região Sudeste, 15,0 na região Sul e 14,9 na região Centro-Oeste.
A nota Máxima no Componente de Formação Geral da prova de alunos de cursos Presenciais de Licenciatura foi 88,1, obtida por, pelo menos, um aluno na região Sudeste.
A menor nota Máxima foi obtida na região Norte ( 72,2 ).
A Mediana do conjunto de Licenciaturas presenciais foi 39,4, sendo a menor Mediana encontrada na região Norte ( 34,4 ), e a maior, encontrada na região Sudeste ( 43,5 ).
A nota Mínima nesta parte foi zero em todas as regiões.
Considerando -se as notas segundo Grande Região, observa -se que não existe diferença estatisticamente significativa entre as Média das notas das Licenciaturas Presenciais, no Componente de Formação Geral, entre as regiões Norte, Nordeste e Centro-Oeste, e entre as regiões Sudeste e Sul.
Somente na região Norte não houve participação de estudantes de cursos a Distância de Licenciatura.
A Média mais alta foi obtida na região Centro-Oeste ( 48,1 ) e a mais baixa foi obtida na região Nordeste ( 29,0 ).
O Desvio padrão das Licenciaturas a Distância foi 15,2, sendo o menor obtido na região Centro-Oeste 13,3 e o maior na região Sudeste ( 15,5 ).
A maior nota Máxima foi 84,8, obtida na região Sudeste, registrada como a nota Máxima de todas as Licenciaturas.
As notas Mínimas foram diferentes de zero nas regiões Nordeste ( 3,2 ) e Centro-Oeste ( 21,0 ).
Na Modalidade de Educação a Distância de Licenciatura, não existe diferenças, estatisticamente significativa ao nível de 95 \%, entre as Médias obtidas nas Grandes Regiões Sudeste e Sul.
Os estudantes de Bacharelado, de todo o Brasil, obtiveram desempenho médio de 37,7 ( Tabela 6.3 ).
Quanto à variabilidade, o Desvio padrão das notas dos estudantes de Bacharelado do Brasil, como um todo, foi 15,3.
A maior Média foi obtida na região Centro-Oeste ( 40,3 ), e a menor, na região Norte ( 23,3 ).
As demais Médias foram : 33,8, na região Nordeste ; 38,7, na região Sudeste ; e 39,3, na região Sul.
Já o maior Desvio padrão foi obtido na região Sudeste ( 15,6 ) e o menor, na região Norte ( 12,2 ).
Na região Nordeste o Desvio padrão foi 14,9, na região Sul foi 14,1 e na região Centro-Oeste foi 15,1.
A maior nota no Componente de Formação Geral da prova do Enade/2021 foi 82,8, obtida por, pelo menos, um estudante na região Sul.
A menor nota Máxima foi obtida na região Norte ( 48,9 ).
A Mediana do Brasil, como um todo, foi 37,5, sendo a menor Mediana encontrada na região Norte ( 24,9 ), e a maior, na região Centro-Oeste ( 41,3 ).
A nota Mínima nessa parte foi igual a zero em três das cinco Grandes Regiões e nas demais foi : Sul ( 6,0 ) e Centro-Oeste ( 3,4 ).
Considerando -se as notas, segundo Grande Região, observa -se que não existe diferença estatisticamente significativa entre as Médias das notas do Componente de Formação Geral obtidas nas regiões Sudeste, Sul e Centro-Oeste.
As diferenças entre as demais Médias são estatisticamente significativas.
As estatísticas das notas dos alunos de cursos Presenciais de Bacharelado das regiões Norte e Nordeste são as mesmas apresentadas para o total de cursos de Bacharelados participantes do Enade em 2021.
A Média mais alta foi obtida nas regiões Sudeste e Sul ( 39,5 ) e a mais baixa foi obtida na região Norte ( 23,3 ).
O Desvio padrão das Licenciaturas Presenciais foi 15,4, sendo o menor obtido na região Norte 12,2 e o maior na região Centro-Oeste ( 15,6 ).
A maior nota Máxima foi 82,8, obtida na região Sul, registrada como a nota Máxima de todas as Licenciaturas.
A nota Mínima nessa parte foi igual a zero em três das cinco Grandes Regiões e nas demais foi : Sul ( 6,0 ) e Centro-Oeste ( 3,4 ).
Considerando -se as notas segundo a Modalidade Presencial e Grande Região, observa -se que não existem diferenças estatisticamente significativas ao nível de 95 \% entre as Médias das regiões do Sudeste, Sul e Centro-Oeste.
Somente nas regiões Norte e Nordeste não foram oferecidos cursos a Distância de Bacharelado A Média mais alta foi obtida na região Centro-Oeste ( 43,3 ) e a mais baixa foi obtida na região Sudeste ( 32,1 ).
O Desvio padrão dos Bacharelados a Distância foi 14,7, sendo o menor obtido na região Centro-Oeste 13,4 e o maior na região Sudeste ( 15,2 ).
A maior nota Máxima foi 73,0, obtida na região Sul, registrada como a nota Máxima de todos os Bacharelados.
As notas Mínimas segundo região, foram : Sudeste ( 4,6 ), Sul ( 8,9 ) e Centro-Oeste ( 18,1 ).
Na Modalidade de Educação a Distância de Bacharelado, não existe diferenças, estatisticamente significativa ao nível de 95 \%, entre as Médias obtidas nas Grandes Regiões Sul e Centro-Oeste.
O Gráfico 6.2 propicia a avaliação do desempenho dos estudantes do Componente de Formação Geral, a partir do histograma da distribuição das notas correspondentes.
A distribuição é unimodal para as notas de Licenciatura e de Bacharelado, com ambas as modas no intervalo ( 30, 40 ].
Na Tabela 6.4, são apresentadas as informações referentes ao desempenho dos concluintes do Componente de Formação Geral, em diferentes agregações : Categoria Administrativa, Organização Acadêmica segundo a Habilitação.
Primeiramente, com respeito aos alunos de Licenciatura, observa -se que não existe diferença estatisticamente significativa entre as Médias dos tipos de Categoria Administrativa : a menor Média obtida por estudantes de Licenciatura de IES Públicas ( 38,3 ) e a maior, pelos de IES Privadas ( 38,9 ).
Nota -se que não há diferença estatisticamente significativa entre as Médias das Universidades e das Faculdades, e das Faculdade e dos Centros Universitários.
As Médias obtidas por Organização Acadêmica foram : 39,0 para as Universidades, 37,3 para os Centros Universitários e 40,0 para as Faculdades.
Dentre os alunos de Bacharelado, observa -se que existe diferença estatisticamente significativa entre as Médias dos tipos de Categoria Administrativa : a maior Média obtida por estudantes de Bacharelado de IES Públicas ( 42,0 ) e a menor, pelos de IES Privadas ( 34,8 ).
Nota -se que não há diferença estatisticamente significativa entre as Médias dos Centros Universitários e das Faculdades.
As Médias foram : 39,8 para as Universidades, 34,7 para os Centros Universitários e 32,3 para as Faculdades, a menor Média.
Em ambas habilitações, a Organização Acadêmica CEFET/IF não teve cursos de Filosofia avaliados. 6.1.3 ESTATÍSTICAS BÁSICAS DO COMPONENTE DE CONHECIMENTO ESPECÍFICO Na Tabela 6.5, são apresentadas as Estatísticas Básicas referentes ao Componente de Conhecimento Específico da Área de Filosofia ( Licenciatura/Bacharelado ) por Grande Região, segundo Habilitação e Modalidade de Ensino.
A nota do componente de Conhecimentos Específicos leva em conta as notas da parte objetiva e a Média das notas obtidas nas três questões discursivas, parte discursiva da Área.
A Média do desempenho dos estudantes de Licenciatura do Brasil, como um todo, foi 52,2.
A maior Média foi obtida na região Sudeste ( 54,4 ), e a menor, na região Norte ( 45,3 ).
As demais Médias foram : 48,9 na região Nordeste, 53,2 na região Sul e 51,1 na região Centro-Oeste.
Quanto à variabilidade das notas, o Desvio padrão do Brasil, como um todo para os alunos de Licenciatura, foi 19,0, sendo o maior Desvio padrão observado na região Centro-Oeste ( 19,9 ) e o menor, na região Norte ( 17,7 ).
Os demais desvios foram : 17,9 na região Nordeste, 19,0 na região Sudeste e 19,2 na região Sul.
A Mediana das notas dos estudantes de Licenciatura de todo o Brasil foi 53,3.
A maior Mediana ocorreu na região Sudeste ( 56,2 ), e a menor, na região Norte ( 45,9 ).
As demais medianas foram : 48,9 na região Nordeste, 54,5 na região Sul e 53,3 na região Centro-Oeste.
A nota Máxima do Brasil, como um todo, foi 97,8, sendo obtida por pelo menos um estudante de Licenciatura na região Sul.
As demais notas Máximas foram : 83,0 na região Norte e 96,3 nas regiões Nordeste, Sudeste e Centro-Oeste.
A nota Mínima foi zero em todas as regiões, com exceção da região Nordeste, onde foi 7,4.
Observa -se que não existe diferença estatisticamente significativa entre as regiões Norte, Nordeste e Centro-Oeste na análise das Médias das notas do Componente de Conhecimento Específico, assim como entre as regiões Sudeste, Sul e Centro-Oeste.
Já para os alunos dos cursos Presenciais de Licenciatura, a Média do desempenho de Conhecimento Específico foi 55,3.
A maior Média foi obtida na região Sudeste ( 60,9 ), e a menor, na região Norte ( 45,3 ).
As demais Médias foram : 51,9 na região Nordeste, 59,6 na região Sul e 50,4 na região Centro-Oeste.
Quanto à variabilidade das notas, o Desvio padrão do conjunto das Licenciaturas presenciais foi 18,4, sendo o maior Desvio padrão observado na região Centro-Oeste ( 20,4 ), e o menor, na região Nordeste ( 17,0 ).
Os demais desvios foram : 17,7 na região Norte, 17,4 na região Sudeste e 18,0 na região Sul.
A Mediana das notas dos estudantes de cursos Presenciais de Licenciatura de todo o Brasil foi 56,3.
A maior Mediana ocorreu na região Sudeste ( 63,6 ), e a menor, na região Norte ( 45,9 ).
As demais medianas foram : 51,9 na região Nordeste, 62,8 na região Sul e 51,8 na região Centro-Oeste.
A nota Máxima das Licenciaturas Presenciais em Filosofia foi 96,3, sendo obtida por, pelo menos, um aluno na região Sudeste, um aluno na região Sul e um aluno na região Centro-Oeste.
As demais notas Máximas foram : 83,0 na região Norte, e 95,5 na região Nordeste.
A nota Mínima foi diferente de zero nas regiões Nordeste ( 7,4 ) e Sul ( 11,1 ).
Considerando -se as Licenciaturas oferecidas presencialmente, observa -se que não existe diferença estatisticamente significativa entre as Médias das notas do Componente de Conhecimento Específico das regiões Sudeste e Sul, das regiões Norte e Centro-Oeste, e das regiões Nordeste e Centro-Oeste.
Em relação aos cursos a Distância de Licenciatura, a Média das notas da prova dos alunos foi 48,7, mais baixa do que os alunos de cursos Presenciais ( 55,3 ).
Somente na região Norte não foram oferecidos cursos a Distância de Licenciatura.
A Média mais alta foi obtida na região Centro-Oeste ( 56,5 ) e a mais baixa foi obtida na região Nordeste ( 35,6 ).
O Desvio padrão das Licenciaturas a Distância foi 19,0, sendo o menor obtido na região Centro-Oeste 14,8 e o maior na região Sul ( 19,0 ).
A maior nota Máxima foi 97,8, obtida na região Sul, registrada como a nota Máxima de todas as Licenciaturas.
As notas Mínimas foram diferentes de zero nas regiões Nordeste ( 7,4 ) e Centro-Oeste ( 18,4 ).
Na Modalidade de Educação a Distância de Licenciatura, não existe diferenças, estatisticamente significativa ao nível de 95 \%, entre as Médias obtidas nas Grandes Regiões Centro-Oeste, Sul e Sudeste.
A Média do desempenho dos estudantes de Bacharelado do Brasil, como um todo, foi 48,9.
A maior Média foi obtida na região Sul ( 52,6 ), e a menor, na região Norte ( 29,6 ).
As demais Médias foram : 42,6, na região Nordeste ; 49,9 na região Sudeste ; e 52,1, na região Centro-Oeste.
Quanto à variabilidade das notas, o Desvio padrão do Brasil, como um todo para os alunos de Bacharelado, foi 18,2, sendo o maior Desvio padrão observado nas região Centro-Oeste ( 18,1 ) e o menor, na região Norte ( 13,1 ).
Os demais desvios foram : 17,6, na região Nordeste, 18,0 na região Sudeste e 16,8, na região Sul.
A Mediana das notas dos estudantes de Bacharelado de todo o Brasil foi 49,0.
A maior Mediana ocorreu na região Sul ( 54,3 ), e a menor, na região Norte ( 29,7 ).
As demais Medianas foram : 40,8 na região Nordeste ; 50,2 na região Sudeste ; e 52,8 na região Centro-Oeste.
A nota Máxima do Brasil, como um todo, foi 91,2, sendo obtida por, pelo menos, um estudante de Bacharelado na região Centro-Oeste.
As demais notas Máximas foram : 60,5, na região Norte ; 83,9 na região Nordeste,90,0 na região Sudeste e 89,3 na região Sul.
A nota Mínima foi zero nas regiões Norte e Sudeste.
A maior nota Mínima foi obtida na região Centro-Oeste ( 14,6 ).
Em relação às notas do Componente de Conhecimento Específico, observa -se que existe diferença estatisticamente significativa entre as duas menores Médias das regiões Norte e Nordeste, e entre essas e as demais.
No entanto, não existe diferença estatisticamente significativa entre as três Médias mais altas, obtidas nas regiões Sudeste, Sul e Centro-Oeste.
Já para os alunos de cursos Presenciais de Bacharelado, a Média do desempenho de Conhecimento foi ligeiramente menor, 48,8.
A maior e a menor Média são respectivamente na Sul ( 53,0 ) e região Norte ( 29,6 ).
As demais Médias foram : 42,6 na região Nordeste, 50,8 na região Sudeste e 51,5 na região Centro-Oeste.
O Desvio padrão dos Bacharelados Presenciais foi 18,3, sendo o menor obtido na região Norte 13,1 e o maior na região Sudeste ( 17,9 ).
As notas Máximas e as notas Mínimas obtidas em todas as regiões para os Bacharelados Presenciais foram as mesmas apresentadas para os Bacharelados do Brasil, como um todo.
Na Modalidade de Educação Presencial de Bacharelado, não existe diferenças, estatisticamente significativa ao nível de 95 \%, entre as Médias obtidas nas Grandes Regiões Sudeste, Sul e Centro-Oeste.
Somente nas regiões Norte e Nordeste não foram oferecidos cursos a Distância de Bacharelado A Média mais alta foi obtida na região Centro-Oeste ( 53,8 ) e a mais baixa foi obtida na região Sudeste ( 43,4 ).
O Desvio padrão dos Bacharelados a Distância foi 17,4, sendo o menor obtido na região Sul ( 15,2 ) e o maior na região Centro-Oeste ( 19,1 ).
A maior nota Máxima foi 88,0, obtida nas regiões Sudeste e Sul, registrada como a nota Máxima de todos os Bacharelados.
As notas Mínimas segundo região, foram : Sudeste ( 9,8 ), Sul ( 13,1 ) e Centro-Oeste ( 22,9 ).
Na Modalidade de Educação a Distância de Bacharelado, não existe diferenças, estatisticamente significativa ao nível de 95 \%, entre as Médias obtidas nas Grandes Regiões Sul e Centro-Oeste.
Assim como nos Gráficos 6.1 e 6.2, no Gráfico 6.3, apresentado a seguir, é possível ser feita uma avaliação do desempenho de concluintes de Filosofia, segundo a Habilitação, em relação ao Componente de Conhecimento Específico, com um histograma da distribuição das notas correspondentes.
Esta também é uma distribuição unimodal para ambas as Habilitações, e o grupo modal é o ( 50 ; 60 ], o mesmo do grupo modal da prova, como um todo ( Gráfico 6.1 ).
Quando comparado o grupo modal da Componente de Conhecimento Específico ( Gráfico 6.3 ) com grupo modal da Formação Geral ( Gráfico 6.2 ), corresponde a dois intervalos acima.
Na Tabela 6.6, é apresentada uma comparação dos resultados em relação a Categoria Administrativa, a Organização Acadêmica e segundo a Habilitação, agora levando -se em conta o desempenho de estudantes do Componente de Conhecimento Específico da prova.
No que se refere a Organização Acadêmica, para os alunos de Licenciatura, a maior Média foi a das Universidades ( 52,8 ), vindo a seguir a das Faculdades ( 51,2 ) e, por fim, o dos Centros Universitários ( 50,8 ).
O maior Desvio padrão foi o dos Centros Universitários ( 19,1 ).
Os Centros Universitários obtiveram a maior nota Máxima ( 97,8 ).
As Universidades obtiveram nota Máxima 96,3, e as Faculdades, 84,5.
As Medianas foram : 54,1 nas Universidades, 51,8 nos Centros Universitários e 52,5 nas Faculdades.
A nota Mínima foi zero para quase todas as Organizações Acadêmicas, exceto para as Faculdades ( 15,5 ).
Observa -se que não existe diferença estatisticamente significativa, ao nível de 95 \%, no Componente de Conhecimento Específico entre as notas das Universidades e Faculdades, e entre as notas dos Centros Universitários e Faculdades.
Quanto a Categoria Administrativa, observa -se um comportamento semelhante àquele da parte de Formação Geral e à prova, como um todo, ou seja, não existe diferença estatisticamente significativa entre as Médias das IES Públicas ( 52,9 ) e as das IES Privadas ( 51,4 ).
No que se refere a Organização Acadêmica, para os alunos de Bacharelado, a maior Média foi a dos Universidade ( 51,2 ), vindo a seguir a dos Centros Universitários ( 46,4 ) e, por fim, a das Faculdades ( 42,4 ).
O maior Desvio padrão, igual ao valor para o Brasil, como um todo, foi o das Universidades ( 18,2 ).
As Universidades obtiveram a maior nota Máxima ( 91,2 ).
Os Centros Universitários obtiveram nota Máxima 88,0 ; as Faculdades, 86,3, a menor delas.
As Medianas foram : 52,8 nas Universidades, 46,5 nos Centros Universitários, 40,8 nas Faculdades.
A nota Mínima foi diferente de zero para todas as Organizações Acadêmicas, exceto para os Centro Universitários, 9,8.
Observa -se que não existe diferença estatisticamente significativa, ao nível de 95 \%, no Componente de Conhecimento Específico entre as Médias dos Centros Universitários e das Faculdades.
Quanto a Categoria Administrativa, observa -se um comportamento semelhante àquele da parte de Formação Geral e à prova, como um todo, ou seja, existe diferença estatisticamente significativa entre as Médias das IES Públicas ( 54,5 ) e as das IES Privadas ( 45,1 ).
Neste caso, também, a maior Média foi obtida por estudantes de IES Públicas de ensino.
Existe diferença estatisticamente significativa, ao nível de 95 \%, na nota do Componente de Conhecimento Específico entre as Categorias Administrativas. 6.2 ANÁLISE DAS QUESTÕES OBJETIVAS Esta seção apresenta estatísticas selecionadas e histogramas das Questões Objetivas de Formação Geral para Licenciatura ( 6.2.1 ) e para Bacharelado ( 6.2.2 ) e de Conhecimento Específico para Licenciatura ( 6.2.3 ) e para Bacharelado ( 6.2.4 ).
São, também, apresentadas e comparadas as Médias das subpopulações caracterizadas por Grande Região. 21 6.2.1 COMPONENTE DE FORMAÇÃO GERAL – FILOSOFIA ( LICENCIATURA ) Na Tabela 6.7a, são apresentadas as Estatísticas Básicas relativas às oito questões objetivas do componente da prova que abrange a Formação Geral dos estudantes de Licenciatura.
A Média do Brasil foi 43,7.
A menor Média foi encontrada na região Norte ( 38,5 ), e a maior, na região Sudeste ( 46,2 ).
As demais Médias foram : 40,1 na região Nordeste, 44,0 na região Sul e 43,0 na região Centro-Oeste.
O Desvio padrão do Brasil foi 20,2, sendo o maior Desvio padrão encontrado na região Norte ( 21,2 ), e o menor Desvio padrão encontrado na região Nordeste ( 19,2 ).
As Medianas do Brasil, como um todo, e de três regiões ( Sudeste, Sul e Centro-Oeste ) foram iguais a 50,0.
Nas regiões Norte e Nordeste as Medianas foram iguais a 37,5.
A nota Máxima 100,0 foi alcançada nas cinco Grandes Regiões.
As notas Mínimas ( 0,0 ) foram iguais para todas as regiões.
Na Tabela 6.8a, são apresentados o Índice de Facilidade e o Índice de Discriminação ( Ponto-Bisserial ) para cada uma das questões objetivas do Componente de Formação Geral.
Quanto ao Índice de Facilidade, foram usadas as seguintes cores para diferenciar o nível de dificuldade da questão : Quanto ao Índice de Discriminação, foram usadas as seguintes cores para qualificar a questão : As questões objetivas do Componente de Formação Geral, segundo o Índice de Facilidade, foram assim avaliadas : Nenhuma teve o Índice de Facilidade classificado como Muito fácil, e uma questão foi tida como Fácil, por ter índice de acertos 0,64.
Quatro questões foram consideradas com índice de dificuldade Médio, situando -se no intervalo entre 0,46 e 0,57 do Índice de Facilidade, ou seja, houve entre 46,0 \% e 57,0 \% de acertos, enquanto três questões apresentaram acertos na faixa entre 18 \% e 30 \%, sendo classificadas como Difícil.
O Índice de Facilidade variou de 0,18 a 0,64, e o de Discriminação, de 0,28 a 0,48.
Seis questões, quanto ao Índice de Discriminação, foram tidas como Muito bom, uma questão foi tida como Bom e uma questão foi tida como Médio.
Na Tabela 6.9a, é apresentada a distribuição das questões, levando -se em conta, simultaneamente, a classificação dos dois índices.
As seis questões que tiveram Índice de Discriminação Muito bom figuraram em três níveis de dificuldade Fácil, Médio e Difícil : uma classificada na categoria Fácil ( a Questão 7 ) do Índice de Facilidade, quatro na categoria Azul para as questões classificadas com índice Muito fácil ( \&gt ; =0,86 ), verde para as questões classificadas com índice Fácil ( 0,61 a 0,85 ), amarelo para as questões classificadas com Médio ( 0,41 a 0,60 ), vermelho para as questões classificadas com Difícil ( 0,16 a 0,40 ) e roxo para as questões classificadas com Muito difícil ( \&lt ; =0,15 ).
Já quanto ao Índice de Discriminação, foram usadas as seguintes cores para qualificar a questão : As questões classificadas com índice Fraco receberam a cor vermelha ( \&lt ; =0,19 ), as classificadas com Médio receberam a cor amarela ( 0,20 a 0,29 ), as classificadas com Bom receberam a cor verde ( 0,30 a 0,39 ) e as classificadas com Muito bom ( \&gt ; =0,40 ) receberam a cor azul.
As questões objetivas do Componente de Formação Geral, segundo o Índice de Facilidade, foram assim avaliadas : Nenhuma teve o Índice de Facilidade classificado como Muito fácil, e uma questão foi tida como Fácil, por ter índice de acertos 0,64.
Quatro questões foram consideradas com índice de dificuldade Médio, situando -se no intervalo entre 0,46 e 0,57 do Índice de Facilidade, ou seja, houve entre 46,0 \% e 57,0 \% de acertos, enquanto três questões apresentaram acertos na faixa entre 18 \% e 30 \%, sendo classificadas como Difícil.
O Índice de Facilidade variou de 0,18 a 0,64, e o de Discriminação, de 0,28 a 0,48.
Seis questões, quanto ao Índice de Discriminação, foram tidas como Muito bom, uma questão foi tida como Bom e uma questão foi tida como Médio.
Médio ( Questões 1, 3, 5 e 6 ), e uma na categoria Difícil ( Questão 4 ).
Em particular, a Questão 3 foi a que apresentou o maior poder discriminatório, com índice 0,48, e foi considerada Médio em termos de Facilidade, com uma proporção de 0,56 acertos.
O máximo de acertos foi alcançado pela Questão 7, com um Índice de Facilidade de 0,64.
A questão 8 foi considerada Difícil, com Índice de Facilidade 0,18.
O Índice de Discriminação foi Médio para a Questão 8 ( 0,28 ).
Na Tabela 6.9a, é apresentada a distribuição das questões, levando -se em conta, simultaneamente, a classificação dos dois índices.
As seis questões que tiveram Índice de Discriminação Muito bom figuraram em três níveis de dificuldade Fácil, Médio e Difícil : uma classificada na categoria Fácil ( a Questão 7 ) do Índice de Facilidade, quatro na categoria Médio ( Questões 1, 3, 5 e 6 ), e uma na categoria Difícil ( Questão 4 ).
Em particular, a Questão 3 foi a que apresentou o maior poder discriminatório, com índice 0,48, e foi considerada Médio em termos de Facilidade, com uma proporção de 0,56 acertos.
O máximo de acertos foi alcançado pela Questão 7, com um Índice de Facilidade de 0,64.
A questão 8 foi considerada Difícil, com Índice de Facilidade 0,18.
O Índice de Discriminação foi Médio para a Questão 8 ( 0,28 ).
No Gráfico 6.4a, para exemplificar, analisa -se o comportamento da Questão de número 3 de Formação Geral.
Trata -se de uma questão considerada com nível Médio de dificuldade, em relação à Facilidade, e a que obteve o maior Índice de Discriminação dessa parte da prova.
Neste gráfico, cada uma das cinco curvas representa o percentual de respostas em determinada alternativa da questão, em função do número de acertos dos estudantes nessa parte da prova ( Formação Geral/Múltipla Escolha ), antes de possíveis eliminações pelo critério Ponto-Bisserial.
A curva em verde corresponde à alternativa B, a alternativa correta para essa questão.
Observa -se que entre os estudantes com menor número de acertos, nessa parte do exame, a situação mais frequente foi a escolha da alternativa E ( em vermelho ), incorreta.
Por exemplo, entre os estudantes que acertaram duas respostas, 41,3 \% escolheram a alternativa E ( em vermelho ), 29,9 \% escolheram a alternativa B ( em verde, gabarito ), 11,8 \% escolheram a alternativa A ( em azul ), 10,2 \% escolheram a alternativa D ( em laranja ) e 6,4 \% escolheram a alternativa C ( em preto ).
Dentre os que acertaram duas respostas entre as questões de múltipla escolha de Formação Geral, 0,0 \% deixou a questão em branco e 0,3 \% marcou mais de uma alternativa, invalidando a questão.
À medida que o número de acertos aumenta, indicando desempenho melhor nessa parte da prova, aumenta, concomitantemente, a proporção de estudantes que selecionaram a alternativa correta B, atingindo 100 \% para os estudantes com oito acertos.
Essa análise permite verificar como a questão discriminou os grupos de desempenho, justificando -se o alto índice ( 0,48 ) obtido na questão.
Cumpre notar que não é possível inferir deste gráfico nem o Índice de Facilidade, que seria uma média da proporção ponderada pela quantidade de estudantes com cada uma das notas, nem o Índice de Discriminação Ponto-Bisserial, por razão equivalente.
No caso extremo, no qual a grande concentração dos acertos dos estudantes fosse abaixo de quatro, o Índice de Facilidade seria obrigatoriamente abaixo de 70,0 \% ( neste exemplo ).
Caso a concentração fosse em seis acertos ou mais, o índice seria obrigatoriamente acima de 80,0 \%.
Os gráficos relativos às demais questões de Formação Geral constam do Anexo IA. 6.2.2 COMPONENTE DE FORMAÇÃO GERAL – FILOSOFIA ( BACHARELADO ) Na Tabela 6.7b, são apresentadas as Estatísticas Básicas relativas às oito questões objetivas válidas do componente da prova que abrange a Formação Geral dos estudantes de Bacharelado.
A Média do Brasil foi 45,9.
A menor Média foi encontrada na região Norte ( 28,0 ), e a maior, na região Centro-Oeste ( 49,8 ).
As demais Médias foram : 42,1, na região Nordeste, 46,9, na região Sudeste e 47,1, na região Sul.
O Desvio padrão do Brasil foi 20,8, sendo o maior Desvio padrão encontrado na região Sudeste ( 21,1 ), e o menor, na região Norte ( 17,8 ).
Os demais desvios foram : 20,2, na região Nordeste, 20,0, nas regiões Sul e Centro-Oeste.
As Medianas foram : 50,0 para as regiões Sudeste, Sul e Centro-Oeste, 37,5 para a região Nordeste e 31,3 para a região Norte.
A nota Máxima 100,0 foi alcançada em três das cinco Grandes Regiões e nas demais foram : 87,5 na região Nordeste e 62,5 na região Norte.
As notas Mínimas ( 0,0 ) foram iguais para todas as regiões.
Na Tabela 6.8b, são apresentados o Índice de Facilidade e o Índice de Discriminação ( Ponto-Bisserial ) para cada uma das questões objetivas do Componente de Formação Geral.
A convenção das cores é a mesma da Tabela 6.8a.
As questões objetivas do Componente de Formação Geral, segundo o Índice de Facilidade, foram assim avaliadas : das oito questões válidas, nenhuma teve o Índice de Facilidade classificado como Muito fácil, duas questões foram tidas como Fácil, três questões foram consideradas com índice de dificuldade Médio e três questões foram consideradas com índice de dificuldade Difícil.
Nenhuma questão apresentou menos de 16 \% de acertos, razão pela qual não houve questão classificada como Muito difícil.
O Índice de Facilidade variou de 0,19 a 0,65 e o de Discriminação, de 0,28 a 0,51.
Sete das oito questões foram tidas como Muito bom e a outra como Médio.
Na Tabela 6.9b, é apresentada a distribuição das questões, levando -se em conta, simultaneamente, a classificação dos dois índices.
As sete questões que tiveram Índice de Discriminação Muito bom figuraram entre três níveis de dificuldade Fácil, Médio e Difícil : duas classificadas na categoria Fácil ( as Questões 3 e 7 ) do Índice de Facilidade, três na categoria Médio ( Questões 1, 5 e 6 ) e duas na categoria Difícil ( Questões 2 e 4 ).
Em particular, a Questão 3 foi a que apresentou o maior poder discriminatório, com índice 0,51, e foi considerada Fácil em termos de facilidade, com uma proporção de 0,62 acertos.
O máximo de acertos foi alcançado pela Questão 7, com um Índice de Facilidade de 0,65.
A Questão 8 foi considerada Difícil, com Índice de Facilidade 0,19.
Seu Índice de Discriminação foi o mais baixo dentre as oito questões válidas : Médio ( 0,28 ).
No Gráfico 6.4b, para exemplificar, analisa -se o comportamento da Questão de número 8 de Formação Geral.
Trata -se de uma questão considerada Fácil, em relação à facilidade e a que obteve o segundo maior Índice de Discriminação dessa parte da prova.
Neste gráfico, cada uma das cinco curvas representa o percentual de respostas em determinada alternativa da questão, em função do número de acertos dos estudantes nessa parte da prova ( Formação Geral/Múltipla Escolha ), antes de possíveis eliminações pelo Na Tabela 6.8b, são apresentados o Índice de Facilidade e o Índice de Discriminação ( Ponto-Bisserial ) para cada uma das questões objetivas do Componente de Formação Geral.
A convenção das cores é a mesma da Tabela 6.8a.
As questões objetivas do Componente de Formação Geral, segundo o Índice de Facilidade, foram assim avaliadas : das oito questões válidas, nenhuma teve o Índice de Facilidade classificado como Muito fácil, duas questões foram tidas como Fácil, três questões foram consideradas com índice de dificuldade Médio e três questões foram consideradas com índice de dificuldade Difícil.
Nenhuma questão apresentou menos de 16 \% de acertos, razão pela qual não houve questão classificada como Muito difícil.
O Índice de Facilidade variou de 0,19 a 0,65 e o de Discriminação, de 0,28 a 0,51.
Sete das oito questões foram tidas como Muito bom e a outra como Médio.
Na Tabela 6.9b, é apresentada a distribuição das questões, levando -se em conta, simultaneamente, a classificação dos dois índices.
As sete questões que tiveram Índice de Discriminação Muito bom figuraram entre três níveis de dificuldade Fácil, Médio e Difícil : duas classificadas na categoria Fácil ( as Questões 3 e 7 ) do Índice de Facilidade, três na categoria Médio ( Questões 1, 5 e 6 ) e duas na categoria Difícil ( Questões 2 e 4 ).
Em particular, a Questão 3 foi a que apresentou o maior poder discriminatório, com índice 0,51, e foi considerada Fácil em termos de facilidade, com uma proporção de 0,62 acertos.
O máximo de acertos foi alcançado pela Questão 7, com um Índice de Facilidade de 0,65.
A Questão 8 foi considerada Difícil, com Índice de Facilidade 0,19.
Seu Índice de Discriminação foi o mais baixo dentre as oito questões válidas : Médio ( 0,28 ).
No Gráfico 6.4b, para exemplificar, analisa -se o comportamento da Questão de número 8 de Formação Geral.
Trata -se de uma questão considerada Fácil, em relação à facilidade e a que obteve o segundo maior Índice de Discriminação dessa parte da prova.
Neste gráfico, cada uma das cinco curvas representa o percentual de respostas em determinada alternativa da questão, em função do número de acertos dos estudantes nessa parte da prova ( Formação Geral/Múltipla Escolha ), antes de possíveis eliminações pelo critério Ponto-Bisserial.
A curva em verde corresponde à alternativa B, a alternativa correta para essa questão.
Observa -se que, entre os estudantes com menor número de acertos, nessa parte do exame, a situação mais frequente foi a escolha da alternativa C ( em preto ), incorreta.
Por exemplo, entre os estudantes que acertaram duas questões, 37,9 \% escolheram a alternativa E ( em vermelho ), 31,7 \% escolheram a alternativa B ( em verde, gabarito ), 11,9 \% escolheram a alternativa A ( em azul ), 10,1 \% escolheram a alternativa D ( em laranja ), e 7,0 \%, a C ( em preto ).
Dentre os que acertaram duas respostas entre as questões de múltipla escolha de Formação Geral, 1,3 \% deixou a questão em branco e 0,0 \% marcou mais de uma alternativa, invalidando a questão.
À medida que o número de acertos aumenta, indicando desempenho melhor nessa parte da prova, aumenta, concomitantemente, a proporção de estudantes que selecionaram a alternativa correta B, atingindo 100 \% para os estudantes com oito acertos.
Essa análise permite verificar como a questão discriminou os grupos de desempenho, justificando -se o alto índice ( 0,51 ) obtido na questão.
Cumpre notar que não é possível inferir deste gráfico nem o Índice de Facilidade, que seria uma média da proporção ponderada pela quantidade de estudantes com cada uma das notas, nem o Índice de Discriminação Ponto-Bisserial, por razão equivalente.
No caso extremo, no qual a grande concentração dos acertos dos estudantes fosse abaixo de quatro, o Índice de Facilidade seria obrigatoriamente abaixo de 70,0 \% ( neste exemplo ).
Caso a concentração fosse em seis acertos ou mais, o índice seria obrigatoriamente acima de 90,0 \%.
Os gráficos relativos às demais questões de Formação Geral constam do Anexo IB. 6.2.3 COMPONENTE DE CONHECIMENTO ESPECÍFICO – FILOSOFIA ( LICENCIATURA ) Na Tabela 6.10a, são apresentadas as Estatísticas Básicas em relação às questões objetivas do Componente de Conhecimento Específico da prova de Filosofia por Grande Região.
Nessa parte da prova, como será discutido ainda nesta seção, nenhuma questão objetiva do Componente de Conhecimentos Específicos foi anulada pela CAA e quatro questões não foram usadas no cômputo das notas por terem sido descartadas pelo critério do Ponto-Bisserial.
Assim, as notas foram calculadas com base em 23 das 27 questões objetivas válidas de Conhecimentos Específicos.
A Média do Brasil deste componente foi 56,0.
A menor Média foi observada na região Norte ( 48,5 ), e a maior, na região Sudeste ( 58,3 ).
O Desvio padrão de todo o Brasil foi 20,0, sendo o menor Desvio padrão encontrado na região Sudeste ( 18,8 ), e o maior, na região Centro-Oeste ( 20,7 ).
A Mediana de todo o Brasil foi 56,5, o mesmo valor da Mediana encontrada na região Centro-Oeste.
Nas regiões Nordeste ( 52,2 ) e Norte ( 47,8 ), os valores foram menores, e, nas regiões Sudeste e Sul ( 60,9 ), o valor foi maior.
A nota Máxima da prova foi 100,00 para as regiões Nordeste, Sudeste e Sul.
Nas demais foram : 91,3 na região Norte e 95,7 na região Centro-Oeste.
Na Tabela 6.11a, são apresentados os Índices de Facilidade e de Discriminação ( Ponto-Bisserial ) das questões objetivas do Componente de Conhecimento Específico, para os estudantes de Filosofia de Licenciatura.
Para facilitar a diferenciação das questões, serão usadas as mesmas cores da Tabela 6.8a e 6.8b para as diferentes classificações dos Índices de Facilidade e de Discriminação.
Observa -se que nenhuma questão foi anulada pela CAA, como já mencionado.
Sendo assim, as análises a seguir dizem respeito às 27 questões válidas.
A partir do Índice de Facilidade obtido, pode -se concluir que quase metade das questões objetivas da prova foram consideradas Médio : das 27 questões válidas, 12 foram classificadas como Médio.
Oito questões foram classificadas como Fácil, e outras sete consideradas como Difícil.
Nenhuma questão foi classificada como Muito fácil ou Muito difícil.
Já quanto aos Índices de Discriminação das questões objetivas do Componente de Conhecimento Específico da prova, obtém -se como resultado a seguinte classificação : 15 das 27 questões válidas foram consideradas Muito Bom, sete delas tiveram Índice de Discriminação Bom, uma teve o Índice de Discriminação Médio e outras quatro tiveram Índice de Discriminação Difícil.
Assim, para 22 em 27 questões, os Índices de Discriminação foram Bom ou Muito bom.
Constata -se, assim, que a prova – no que se refere ao Componente de Conhecimento Específico – possuía boa capacidade de discriminar entre aqueles que dominam ou não o conteúdo.
A nota Mínima foi zero, obtida em todas as regiões, com exceção da região Nordeste onde foi 8,7.
Na Tabela 6.11a, são apresentados os Índices de Facilidade e de Discriminação ( Ponto-Bisserial ) das questões objetivas do Componente de Conhecimento Específico, para os estudantes de Filosofia de Licenciatura.
Para facilitar a diferenciação das questões, serão usadas as mesmas cores da Tabela 6.8a e 6.8b para as diferentes classificações dos Índices de Facilidade e de Discriminação.
Observa -se que nenhuma questão foi anulada pela CAA, como já mencionado.
Sendo assim, as análises a seguir dizem respeito às 27 questões válidas.
A partir do Índice de Facilidade obtido, pode -se concluir que quase metade das questões objetivas da prova foram consideradas Médio : das 27 questões válidas, 12 foram classificadas como Médio.
Oito questões foram classificadas como Fácil, e outras sete consideradas como Difícil.
Nenhuma questão foi classificada como Muito fácil ou Muito difícil.
Já quanto aos Índices de Discriminação das questões objetivas do Componente de Conhecimento Específico da prova, obtém -se como resultado a seguinte classificação : 15 das 27 questões válidas foram consideradas Muito Bom, sete delas tiveram Índice de Discriminação Bom, uma teve o Índice de Discriminação Médio e outras quatro tiveram Índice de Discriminação Difícil.
Assim, para 22 em 27 questões, os Índices de Discriminação foram Bom ou Muito bom.
Constata -se, assim, que a prova – no que se refere ao Componente de Conhecimento Específico – possuía boa capacidade de discriminar entre aqueles que dominam ou não o conteúdo.
O Índice de Facilidade variou de 0,20 a 0,79 e o de Discriminação, de 0,02 a 0,53.
Na Tabela 6.12a, é apresentada a distribuição das questões levando -se em conta, simultaneamente, a classificação dos dois índices.
Dentre as questões que alcançaram os maiores Índices de Discriminação, 15 delas foram classificadas com Índice de Discriminação Muito bom.
As questões 11, 14, 22, 27, 28 e 32 apresentam, respectivamente, os seguintes Índices de Facilidade : 0,62, 0,72,079, 0,71, 0,72 e 0,69, sendo classificadas como Fácil.
Já as questões 17, 20, 21, 23, 29, 34 e 35, apresentam os seguintes Índices de Facilidade : 0,58, 0,60, 0,54, 0,56, 0,58, 0,47 e 0,56, sendo classificas como Médio.
Por fim, as questões 18 e 31 apresentam os seguintes Índices de Facilidade : 0,28 e 0,39 que receberam a classificação Difícil.
A Questão de número 15 foi a mais difícil entre as 27 questões específicas, com baixo Índice de Facilidade, apenas 20 \% de acertos.
Essa questão apresentou poder discriminatório muito baixo, 0,02, o que comprova ter sido a mais difícil para os estudantes.
Destaca -se, também, as questões 13, 24 e 30 com os respectivos Índices de Facilidade iguais a 0,23, 0,28 e 0,33, o que, em termos percentuais, corresponde a 23,0 \%, 28,0 \% e 33,0 \% de estudantes que responderam acertadamente, sendo, também, 0,10, 0,09 e 0,08 seus Índices de Discriminação.
Tais questões foram, portanto, pelo critério Ponto-Bisserial, consideradas inadequadas.
Por isso, as questões 13, 15, 24 e 30 foram eliminadas do cômputo da nota final.
A título de exemplo das análises do comportamento das questões objetivas, no Gráfico 6.5a, analisa -se a Questão 27 do Componente de Conhecimento Específico.
Essa questão foi considerada, pelas respostas dos estudantes avaliados na prova, como uma questão com nível de Facilidade Fácil, 0,71, ou seja, 71 \% dos estudantes assinalaram, acertadamente, a opção A correspondente ao gabarito.
Seu Índice de Discriminação foi igual a 0,53, classificado como Muito bom, o maior valor de Discriminação.
Neste gráfico, cada uma das cinco curvas representa o percentual de respostas em determinada alternativa da Questão 27, em função do número de acertos dos estudantes nessa parte da prova, antes de possíveis eliminações de questões pelo critério Ponto-Bisserial.
A alternativa correta B, representada no gráfico pela curva em verde, foi escolhida em maiores proporções pelos estudantes com desempenho melhor nessa parte da prova.
Já as alternativas incorretas, também denominadas distratores, foram selecionadas, principalmente, por aqueles com notas mais baixas.
No caso de Filosofia ( Licenciatura ), como se observa no eixo horizontal do Gráfico 6.5a, nenhum estudante acertou mais do que 26 questões dentre as 27 questões objetivas válidas, considerando -se o gabarito original.
A eliminação das quatro questões com Índice de Discriminação Fraco do cômputo da nota é que possibilitou a existência de nota Máxima igual a 100,0 ( visto na Tabela 6.10a ).
Observa -se que a soma não é 100 \%, por causa das questões não respondidas ou com mais de uma opção marcada.
Aqueles com nenhum acerto, na sua totalidade, marcaram a opção A ( linha azul ).
A proporção de estudantes que selecionaram a resposta correta B aumenta gradativamente, chegando a atingir 100 \% para 23 acertos ou mais, enquanto a proporção dos que escolheram alternativas incorretas decai, a partir de três acertos, em função do número de acertos nessa parte da prova.
Os gráficos relativos às demais questões do Conhecimento Específico constam do Anexo IA. 6.2.4 COMPONENTE DE CONHECIMENTO ESPECÍFICO – FILOSOFIA ( BACHARELADO ) Na Tabela 6.10b, são apresentadas as Estatísticas Básicas em relação às questões objetivas do Componente de Conhecimento Específico da prova de Filosofia por Grande Região.
Nessa parte da prova, nenhuma questão foi anulada pela Comissão Assessora da Área.
Além disso, como será discutido ainda nesta seção, nenhuma questão objetiva do Componente de Conhecimentos Específicos foi descartada pelo critério do Ponto-Bisserial.
Na Tabela 6.11b, são apresentados os Índices de Facilidade e de Discriminação ( Ponto-Bisserial ) das questões objetivas do Componente de Conhecimento Específico, para os estudantes de Filosofia de Bacharelado.
Para facilitar a diferenciação das questões, serão usadas as mesmas cores da Tabela 6.8a, 6.8b e 6.11a para as diferentes classificações dos Índices de Facilidade e de Discriminação.
Observa -se que nenhuma questão foi anulada pela Comissão Assessora da Área ( CAA ) e que as análises dizem respeito às 27 questões válidas.
A partir do Índice de Facilidade obtido, pode -se concluir que mais da metade das questões objetivas da prova foram considerados Médio, 15 das 27 questões válidas, e Assim, as notas foram calculadas com base em todas 27 questões objetivas válidas de Conhecimentos Específicos.
A Média do Brasil deste componente foi 51,8.
A menor Média foi observada na região Norte ( 32,6 ), e a maior, na região Centro-Oeste ( 56,5 ).
O Desvio padrão de todo o Brasil foi 19,0, sendo o menor Desvio padrão encontrado na região Norte ( 14,1 ), e o maior, na região Centro-Oeste ( 18,9 ).
A Mediana de todo o Brasil foi 53,8, o mesmo valor da Mediana encontrada na região Sudeste.
Nas regiões Norte e Nordeste a Mediana foi menor ( 32,7 e 42,3 ) e nas regiões Sul e Centro-Oeste, foi maior ( 57,7, em ambas ).
A nota Máxima da prova foi 96,2, obtida nas questões objetivas do Componente de Conhecimento Específico, por, pelo menos, um estudante das regiões Sudeste e Sul.
A nota Mínima foi zero nas regiões Norte e Sudeste.
Na Tabela 6.11b, são apresentados os Índices de Facilidade e de Discriminação ( Ponto-Bisserial ) das questões objetivas do Componente de Conhecimento Específico, para os estudantes de Filosofia de Bacharelado.
Para facilitar a diferenciação das questões, serão usadas as mesmas cores da Tabela 6.8a, 6.8b e 6.11a para as diferentes classificações dos Índices de Facilidade e de Discriminação.
Observa -se que nenhuma questão foi anulada pela Comissão Assessora da Área ( CAA ) e que as análises dizem respeito às 27 questões válidas.
A partir do Índice de Facilidade obtido, pode -se concluir que mais da metade das questões objetivas da prova foram considerados Médio, 15 das 27 questões válidas, e nenhuma como Muito difícil.
Quatro questões foram classificadas como Fácil, e outras oito consideradas como Difícil.
Nenhuma questão foi classificada como Muito fácil.
O Índice de Facilidade variou de 0,24 a 0,83 e o de Discriminação, de 0,05 a 0,56. nenhuma como Muito difícil.
Quatro questões foram classificadas como Fácil, e outras oito consideradas como Difícil.
Nenhuma questão foi classificada como Muito fácil.
Já quanto aos Índices de Discriminação das questões objetivas do Componente de Conhecimento Específico da prova, obtém -se como resultado a seguinte classificação : 14 das 27 questões foram consideradas como Muito Bom, enquanto oito delas tiveram Índice de Discriminação Bom.
Assim, para 22 em 27 questões, os Índices de Discriminação foram Bom ou Muito bom.
Dentre as demais, quatro delas foram classificadas como Médio, e uma como Fraco, sendo cinco, por conseguinte, a quantidade de questões nos dois patamares mais baixos de discriminação.
Constata -se, assim, que a prova – no que se refere ao Componente de Conhecimento Específico – possuía boa capacidade de discriminar entre aqueles que dominam ou não o conteúdo.
O Índice de Facilidade variou de 0,24 a 0,83 e o de Discriminação, de 0,05 a 0,56.
Na Tabela 6.12b, é apresentada a distribuição das questões levando -se em conta, simultaneamente, a classificação dos dois índices.
Dentre as questões que alcançaram os maiores Índices de Discriminação, uma delas obteve o maior deles, a de número 28.
Seu índice foi 0,56, e, quanto ao Índice de Facilidade, foi classificada como Médio, com 53 \% dos estudantes marcando a opção correta.
A Questão de número 15 foi a mais difícil entre as 27 questões específicas, com baixo Índice de Facilidade, apenas 24 \% de acertos.
Essa questão apresentou poder discriminatório muito baixo, 0,05, o que comprova ter sido a mais difícil para os estudantes.
Tal questão foi, portanto, pelo critério Ponto-Bisserial, considerada inadequada.
Por isso, a Questão 15 foi eliminada do cômputo da nota final.
A título de exemplo das análises do comportamento das questões objetivas, no Gráfico 6.5b, analisa -se a Questão 28 do Componente de Conhecimento Específico.
Essa questão foi considerada, pelas respostas dos estudantes avaliados na prova, como uma questão com nível de facilidade Médio, 0,53, ou seja, 53 \% dos estudantes assinalaram, acertadamente, a opção E ( em vermelho ), correspondente ao gabarito.
Seu Índice de Discriminação foi igual a 0,56, classificado como Muito bom, o maior valor de discriminação.
Neste gráfico, cada uma das cinco curvas representa o percentual de respostas em determinada alternativa da Questão 28, em função do número de acertos dos estudantes nessa parte da prova, antes de possíveis eliminações de questões pelo critério Ponto-Bisserial.
A alternativa correta E, representada no gráfico pela curva em vermelho, foi escolhida em maiores proporções pelos estudantes com desempenho melhor nessa parte da prova.
Já as alternativas incorretas, também denominadas distratores, foram selecionadas, principalmente, por aqueles com notas mais baixas.
No caso de Filosofia ( Bacharelado ), como se observa no eixo horizontal do Gráfico 6.5b, nenhum estudante acertou mais do que 26 questões dentre as 27 questões objetivas válidas, considerando -se o gabarito original.
Mesmo com a eliminação de uma questão com Índice de Discriminação Fraco do cômputo da nota, nenhum participante atingiu a nota Máxima, igual a 100,0, nessa parte da prova ( visto na Tabela 6.10b ).
Observa -se que a soma não é 100 \%, por causa das questões não respondidas ou com mais de uma opção marcada.
Aqueles com nenhum acerto, na sua quase totalidade, marcaram a opção C ( em preto ).
A proporção de estudantes que selecionaram a resposta correta E aumenta gradativamente, chegando a atingir 100 \% para 24 acertos ou mais, enquanto a proporção dos que escolheram alternativas incorretas decai, a partir de cinco ou seis acertos, em função do número de acertos nessa parte da prova.
Os gráficos relativos às demais questões do Conhecimento Específico constam do Anexo IB. 6.3 ANÁLISE DAS QUESTÕES DISCURSIVAS Esta seção apresenta estatísticas selecionadas e histogramas das Questões Discursivas de Formação Geral ( 6.3.1 ) e Conhecimento Específico ( 6.3.2 ).
São, também, apresentadas e comparadas as médias de subpopulações, caracterizadas por Grande Região.
Cumpre notar que nem todas as questões passam pelo crivo dos corretores.
As respostas depois de digitalizadas passam por um pré-processamento de identificação automático.
Algumas destas questões são definidas como um “ branco automático ” ( que não são enviadas para os professores para correção, a não ser que na segunda fase, a de inspeção visual, seja constatada uma resposta ).
O critério usado neste pré-processamento, é verificar para cada questão, a partir da quantidade de PIXELS, o que provavelmente, é uma resposta em branco.
A partir da separação destes documentos que foram considerados em branco, eles são enviados para uma verificação visual, por uma equipe de colaboradores, devidamente treinados.
Deste modo, quaisquer erros, deste pré-processamento, são encontrados, e as questões preenchidas, são enviadas para a correção dos professores.
Eventualmente, algumas respostas em branco, não são detectadas neste pré-processamento automático.
A presença de pequenos riscos, ou manchas, podem fazer o programa não identificar uma questão como um possível branco.
Estas são enviadas para a correção, mas são classificadas como “ branco ”, pelos professores quando da correção.
No caso da Área de Filosofia ( Licenciatura ), foram identificados como brancos “ automáticos ”, 1.449 respostas de questões discursivas de Conhecimento Específico e 737 de Formação Geral.
Pelos professores corretores foram identificadas mais 12 e seis questões em branco, respectivamente, para Conhecimento Específico e para Formação Geral.
No caso da Área de Filosofia ( Bacharelado ), foram identificados como brancos “ automáticos ”, 566 respostas de questões discursivas de Conhecimento Específico e 284 de Formação Geral.
Pelos professores corretores foram identificadas mais quatro e duas questões em branco, respectivamente, para Conhecimento Específico e para Formação Geral.
Uma outra situação a ser mencionada é a ocorrências de atendimento a portadores de necessidades especiais.
O Anexo XI apresenta, para o exame como um todo e para a Área de Filosofia ( Licenciatura/Bacharelado ), por tipo de deficiência, o protocolo usado para permitir a correção das respostas dos alunos em cada situação. 6.3.1 COMPONENTE DE FORMAÇÃO GERAL As análises dos resultados de desempenho dos estudantes de Filosofia ( Licenciatura/Bacharelado ), nas duas questões discursivas relativas a Formação Geral, encontram -se na Tabela 6.13 e no Gráfico 6.6.
Aqui analisa -se a média das notas de cada Questão Discursiva de Formação Geral que leva em conta as correções realizadas por duas bancas, uma que analisa e pontua aspectos relativos ao conteúdo das respostas, seguindo o padrão de respostas divulgado pelo Inep ( ver Anexos IXA e IXB ), outra que avalia o desempenho linguístico expresso pela redação das respostas, seguindo um padrão de respostas específico da Área de Língua Portuguesa ( ver Anexos IXA e IXB ).
A nota de cada questão é composta de 80 \% da nota da correção de conteúdo e 20 \% da nota de desempenho linguístico.
Na Tabela 6.13, observa -se que a nota Média nesse conjunto de questões foi inferior à obtida nas objetivas pelos estudantes de Licenciatura.
Os estudantes de Licenciatura de todo o Brasil obtiveram, em Formação Geral, Média 43,7 ( Tabela 6.7a ) nas questões objetivas e 30,9 nas questões discursivas.
Pode -se notar, também, que o Desvio padrão nesse conjunto de questões foi menor do que o obtido nas objetivas : 20,2 nas questões objetivas e 16,8 nas questões discursivas.
A maior Média foi obtida na região Sudeste ( 32,9 ) e a menor, na região Norte ( 26,5 ).
A Mediana dos alunos de Licenciatura de todo o Brasil, neste componente, foi 31,5.
Nas regiões Norte, Nordeste e Sul, a Mediana foi menor ( 27,5, 30,5 e 30,5, respectivamente ), e, nas demais regiões, foi maior que o valor para o Brasil : Sudeste e Centro-Oeste ( 33,0, ambas ).
A nota Máxima ( 89,0 ) foi obtida na região Sudeste, sendo 65,5, a Máxima na região Norte ; 73,5, na região Nordeste ; 81,0, na região Sul ; e 72,0, na região Centro-Oeste.
A nota Mínima ( 0,0 ) foi a mesma em todas as regiões do Brasil.
Para os alunos de cursos Presenciais de Licenciatura, a nota Média, nesse conjunto de questões, foi inferior da obtida nas objetivas.
Estes estudantes obtiveram, em Formação Geral, Média 45,2 ( Tabela 6.7a ) nas questões objetivas e 31,2 nas questões discursivas.
Pode -se notar também que o Desvio padrão nesse conjunto de questões foi um pouco menor do obtido nas objetivas, 19,9 nas questões objetivas e 17,0 nas questões discursivas.
A maior Média foi obtida na região Sudeste ( 33,8 ), e a menor, na região Norte ( 26,5 ).
A Mediana dos alunos de cursos Presenciais de Licenciatura, nas questões discursivas deste componente, foi 31,5.
Nas regiões Norte ( 27,5 ), Nordeste ( 31,0 ), Sul ( 31,0 ) e Centro-Oeste ( 30,0 ) a Mediana foi menor e na região Sudeste ( 33,5 ) foi maior que o valor para o Brasil.
A nota Máxima ( 89,0 ) foi a mesma na região Sudeste, sendo 65,5 na região Norte, 73,5 na região Nordeste, 81,0, na região Sul e 71,5, na região Centro-Oeste.
A nota Mínima ( 0,0 ) foi a mesma em todas as regiões do Brasil.
Também para os alunos de cursos a Distância de Licenciatura, a participação não ocorreu apenas na região Norte.
Observa -se que a nota Média nesse conjunto de questões foi inferior da obtida nas objetivas.
Estes estudantes obtiveram, em Formação Geral, Média 41,9 ( Tabela 6.7a ) nas questões objetivas e 30,6 nas questões discursivas.
Pode -se notar também que o Desvio padrão nesse conjunto de questões foi um pouco menor do obtido nas objetivas, 20,3 nas questões objetivas e 16,5 nas questões discursivas.
A maior Média foi obtida na região Centro-Oeste ( 45,4 ), e a menor, na região Nordeste ( 26,3 ).
A Mediana dos alunos de cursos a Distância de Licenciatura, nas questões discursivas deste componente, foi 31,5.
Na região Nordeste ( 26,0 ) e na região Sul ( 30,5 ) a Mediana foi menor e nas demais regiões foi maior que o valor para o Brasil : Sudeste ( 33,0 ) e Centro-Oeste ( 48,3 ).
A nota Máxima ( 80,0 ) foi a mesma na região Sul, sendo 77,5 na região Sudeste, 72,0, na região Centro-Oeste e 65,0 na região Nordeste.
A nota Mínima ( 0,0 ) foi a mesma em todas as regiões do Brasil.
Considerando -se agora os alunos de cursos de Bacharelado, na Tabela 6.13, observa -se que a nota Média nesse conjunto de questões foi inferior à obtida nas objetivas pelos estudantes de Bacharelado.
Os estudantes de Bacharelado de todo o Brasil obtiveram, em Formação Geral, Média 45,9 ( Tabela 6.7b ), nas questões objetivas e 25,4, nas questões discursivas.
Pode -se notar, também, que o Desvio padrão nesse conjunto de questões foi menor do que o obtido nas objetivas : 20,8, nas questões objetivas e 15,6, nas questões discursivas.
A maior Média foi obtida na região Sul ( 27,6 ), e a menor, na região Norte ( 16,2 ).
A Mediana dos alunos de Bacharelado de todo o Brasil, neste componente, foi 24,5.
A Mediana foi menor nas regiões Norte ( 13,3 ), Nordeste ( 19,5 ), e, nas demais regiões foi maior que o valor para o Brasil : Sudeste ( 25,0 ), Sul ( 27,3 ) e Centro-Oeste ( 25,5 ).
A nota Máxima ( 74,5 ) foi obtida na região Sudeste, sendo 44,5, a nota Máxima na região Norte ; 70,5, na região Nordeste ; 70,0, na região Sul ; e 67,0, na região Centro-Oeste.
A nota Mínima ( 0,0 ) foi a mesma em todas as regiões do Brasil.
Considerando -se agora os alunos de cursos Presenciais de Bacharelado, observa -se que a nota Média nesse conjunto de questões foi inferior da obtida nas objetivas.
Estes estudantes obtiveram, em Formação Geral, Média 46,0 ( Tabela 6.7b ) nas questões objetivas e 25,3 nas questões discursivas.
Pode -se notar também que o Desvio padrão nesse conjunto de questões foi menor do obtido nas objetivas, 20,9 nas questões objetivas e 15,8 nas questões discursivas.
A maior Média foi obtida na região Sul ( 27,5 ), e a menor, na região Norte ( 16,2 ).
Também para os alunos de cursos a Distância de Bacharelado, a participação não ocorreu nas regiões Norte e Nordeste.
Observa -se que a nota Média nesse conjunto de questões foi inferior da obtida nas objetivas.
Estes estudantes obtiveram, em Formação Geral, Média 44,9 ( Tabela 6.7a ) nas questões objetivas e 26,0 nas questões discursivas.
Pode -se notar também que o Desvio padrão nesse conjunto de questões foi um pouco menor do obtido nas objetivas, 20,7 nas questões objetivas e 14,5 nas questões discursivas.
A maior Média foi obtida na região Centro-Oeste ( 15,3 ), e a menor, na região Sudeste ( 12,7 ).
A Mediana dos alunos de cursos a Distância de Bacharelado, nas questões discursivas deste componente, foi 25,5.
Na região Sudeste ( 21,5 ) a Mediana foi menor e nas demais regiões foi maior que o valor para o Brasil : Sul ( 28,0 ) e Centro-Oeste ( 31,8 ).
A nota Máxima ( 70,0 ) foi a mesma na região Sul, sendo 50,5 na região Sudeste e 66,5, na região Centro-Oeste.
A nota Mínima ( 0,0 ) foi a mesma em todas as regiões do Brasil.
O Gráfico 6.6 representa a distribuição das notas nas Questões Discursivas do Componente de Formação Geral.
A moda desta distribuição ocorre no intervalo ( 30 ; 40 ] para os cursos de Licenciatura ( 24,0 \% ) e no intervalo, ( 20 ; 30 ], para o Bacharelado ( 24,9 \% ).
Destaca -se também o intervalo [ 0 ; 10 ] com uma proporção de 18,0 \% do total de notas do Bacharelado e de 12,7 \% do total de notas da Licenciatura.
Destaca -se, também, o intervalo [ 0 ; 10 ], que se inclui, além da nota zero, a frequência de estudantes que deixaram esse tipo de questão em branco.
Na sequência, os resultados verificados para cada uma das questões discursivas de Formação Geral serão apresentados, estabelecendo -se relações com os conteúdos abordados em cada uma delas.
Os comentários das Bancas de docentes corretores, a respeito do observado na correção das respostas dos estudantes, suas impressões e conclusões serão apresentados junto à análise de cada questão.
Cumpre esclarecer que, tendo em vista que as questões discursivas de Formação Geral são padronizadas, ou seja, constam de todas as provas, os comentários da Banca são os mesmos para todas as carreiras acadêmicas, sendo direcionados a todos os estudantes que participaram do Enade/2021.
A seguir, serão analisados os desempenhos linguísticos dos estudantes da Área de Filosofia ( Licenciatura/Bacharelado ) nas duas questões discursivas de Formação Geral do Enade/2021, apresentando -se os resultados obtidos e os comentários da banca de correção para cada questão. 6.3.1.1 ANÁLISE DE CONTEÚDO DA QUESTÃO DISCURSIVA 1 DO COMPONENTE DE FORMAÇÃO GERAL Os dados de Filosofia ( Licenciatura/Bacharelado ), obtidos a partir das respostas à Questão 1, encontram -se na Tabela 6.14 e no Gráfico 6.7.
Nessa questão – de desempenho muito semelhante ao da outra questão de Formação Geral –, os estudantes de Licenciatura de todo o Brasil tiveram Média 22,5.
A maior Média para a Questão 1 foi obtida na região Sul ( 25,5 ), e a menor, na região Nordeste ( 18,0 ).
Quanto à variabilidade das notas, o Desvio padrão de todo o Brasil foi 19,0.
O menor Desvio padrão foi obtido na região Nordeste ( 17,4 ), e o maior Desvio padrão foi obtido pela região Sudeste ( 19,6 ).
A Mediana do Brasil, para os alunos de Licenciatura, foi 25,0, a mesma obtida na região Sudeste.
A regiões Norte e Nordeste obtiveram Mediana menores do que a do Brasil, sendo elas : 22,5 e 15,0, respectivamente.
Já as regiões Sul e Centro-Oeste obtiveram Mediana igual a 30,0.
A nota Máxima da Questão Discursiva 1 foram : 95,0 para as regiões Nordeste, Sudeste e Sul, e 80,0 para as regiões Norte e Centro-Oeste.
Dentre os alunos de Bacharelado, os dados obtidos a partir das respostas à Questão 1 – de desempenho inferior ao verificado na outra questão de Formação Geral – os estudantes de Bacharelado de todo o Brasil, tiveram Média 17,1.
A maior Média para a Questão 1 foi obtida na região Sudeste ( 18,6 ), e a menor, na região Norte ( 8,7 ).
Quanto à variabilidade das notas, o Desvio padrão de todo o Brasil foi 18,0.
O menor Desvio padrão foi obtido na região Norte ( 15,3 ), e o maior Desvio padrão foi obtido pela região Sul ( 18,6 ).
A Mediana do Brasil, para os alunos de Bacharelado, como um todo, foi 15,0, sendo igual a todas as regiões, com exceção da região Norte, onde a Mediana foi igual a zero, significando que pelo menos 50 \% dos participantes zeraram essa questão.
As notas Máxima e Mínima da Questão Discursiva 1 foram as mesmas para todas as regiões do Brasil, respectivamente, 85,0 e zero, com exceção das notas Máximas das regiões Norte e Centro-Oeste, que foram, respectivamente, 55,0 e 70,0.
No Gráfico 6.7, mostra -se a distribuição das notas na Questão Discursiva 1 do Componente de Formação Geral.
Observa -se, para os alunos de Licenciatura, que 9,2 \% dos participantes deixaram em branco a resposta a essa questão.
Somados à frequência dos que, tendo respondido, receberam nota zero, chega a 28,1 \% o contingente daqueles que zeraram questão.
O intervalo ( 20 ; 30 ] caracteriza máximo local com 29,2 \% de frequência, cujo valor é muito próximo do percentual de alunos que zeraram a questão.
Dentre os alunos de Bacharelado, observa -se que 9,4 \% dos participantes deixaram em branco a resposta a essa questão e que a classe modal ocorre na nota zero, daqueles que responderam à questão, com frequência igual a 31,0 \%.
Somadas essas duas frequências, chega 40,4 \% o contingente daqueles que zeraram a questão.
Observa -se, ainda, um máximo local no intervalo ( 10 ; 20 ], com 22,5 \%. 6.3.1.2 COMENTÁRIOS SOBRE A CORREÇÃO DE CONTEÚDO DAS RESPOSTAS À QUESTÃO DISCURSIVA 1 O enunciado da Questão1 continha dois textos e um comando.
O primeiro texto trazia uma abordagem sobre as tensões entre cultura, arte e censura sob o \textbf{recorte} do inciso IX do Art. 5º da Constituição Federal de 1988, que diz ser livre a expressão da atividade intelectual, artística, científica e de comunicação, independentemente de censura ou licença.
O segundo texto apresentava o inciso IX do parágrafo 5º da Constituição Federal, em que é garantida a livre expressão da atividade artística, científica e de comunicação, independentemente de censura ou licença.
O comando solicitava que o respondente discorresse a respeito de dois aspectos.
Primeiramente, o estudante deveria ser capaz de dissertar sobre as relações entre os conceitos de arte, cultura e censura, à luz da ideia de liberdade artística garantida pela Constituição Federal de 1988.
O estudante também deveria ser capaz de apresentar duas ações educativas que pudessem contribuir “ para minimizar essas tensões ” ( tensões entre arte, cultura e censura ) e garantir a liberdade artística prevista pela lei.
A questão, portanto, estava de acordo com os temas previstos nas Diretrizes Curriculares para a Formação Geral : ética, democracia e cidadania ; cultura, arte e comunicação ; educação e desenvolvimento humano e social.
Além disso, o comando permitia avaliar o desenvolvimento das seguintes competências previstas nas Diretrizes Curriculares para a Formação Geral : planejar e elaborar projetos de ação e intervenção a partir da análise de necessidades em contextos diversos ; ler, produzir e interpretar textos com clareza e coerência ; identificar, compreender e analisar situações-problema utilizando pensamento lógico e sistêmico, estabelecendo relações de causalidade.
O enunciado da questão era claro, empregava linguagem adequada e continha informações suficientes para que os egressos construíssem seus argumentos.
Entretanto, apesar dessas características favoráveis, a temática era sofisticada, exigia reflexão e estabelecimento de relações a partir de textos concisos e densos.
Além disso, os conceitos de arte e cultura expressos pelo Texto I excluíam inúmeras outras definições reconhecidas em diferentes áreas das ciências humanas, em alguns casos divergentes daquelas apresentadas.
Os conhecimentos e o nível de profundidade exigidos pelo padrão de resposta eram adequados para egressos de cursos superiores.
Deve -se destacar que o padrão admitia um espectro amplo para a construção da argumentação das respostas, viabilizando a aceitação de outras interpretações dos conceitos de arte e cultura, incluindo possíveis contestações em relação às ideias apresentados no Texto I.
A questão pode ser considerada difícil, já que a primeira parte exigia uma resposta que articulasse quatro conceitos complexos – arte, cultura, censura e liberdade artística –, e a segunda parte exigia o conhecimento do que seja uma ação educativa, além da proposição de ações que objetivamente contribuíssem para minimizar as tensões entre arte, cultura e censura, além de garantir a liberdade artística.
Essas observações, aliadas ao fato do espaço e do tempo limitados para elaboração de uma resposta tão extensa, elevaram consideravelmente o nível de dificuldade da questão.
A correção confirmou a dificuldade dos estudantes para abordar a temática, evidenciou grande diversidade na compreensão dos conceitos envolvidos e falta de clareza do que seja uma ação educativa.
Nas respostas à primeira parte do comando, os estudantes raramente abordaram os quatro conceitos que deveriam estar presentes na resposta.
A maioria se limitou a relacionar arte e cultura com comentários que repetiam trechos dos enunciados, como arte e cultura"andam de mãos dadas".
Assim, apesar de o Texto I sugerir o entendimento de que a cultura é uma esfera reguladora do comportamento e de que a arte é o elemento de desequilíbrio, de contestação, a quase totalidade das respostas não fez uso dessa interpretação, tratando arte e cultura quase como sinônimos.
Apesar de raras, houve respostas nas quais os estudantes defenderam a existência de múltiplas culturas, muitas delas não hegemônicas, tão suscetíveis aos dispositivos de controle quanto a arte.
Houve também respostas que colocavam em questão os critérios, muitas vezes elitistas e eurocêntricos, utilizados para produzir hierarquização entre as artes, chegando a inviabilizar o reconhecimento de outras tantas manifestações artísticas.
Foi pequeno o grupo que, além de relacionar arte e cultura, foi capaz de articulá -las com o conceito de censura trazendo -a, quase sempre, como herança do período militar.
Entretanto, nesse grupo, outras abordagens de censura foram mapeadas no decorrer das correções : cancelamento em redes sociais ; contrainformação ou manipulação dos meios de comunicação social ; além de outros recursos censórios mais indiretos, como os processos judiciais, a suspensão de financiamentos públicos a projetos culturais e a autocensura.
Em relação à segunda parte do comando, era esperado que as/os estudantes egressas/os citassem duas ações educativas sob o \textbf{recorte} de que minimizassem tensões observadas entre arte, cultura e censura e garantissem a liberdade artística prevista pela lei.
Na maioria das respostas, não se verificou essa relação entre a proposição e seu objetivo/resultado ; muitas ações foram apenas citadas sem explicações.
Esse fato reduziu bastante a quantidade de respostas com pontuação integral no referido item.
Importante ressaltar que o padrão de resposta indicou exemplos de ações educativas no campo da educação não formal.
Entretanto, por não ser algo explícito no comando, foram também aceitas ações referentes ao campo da educação formal.
O mais recorrente nas respostas foi a indicação de ações educativas em sala de aula e visitas a museus e a exposições, mediadas pelo próprio professor da turma.
Entretanto, houve respostas em que os estudantes ampliaram tanto o espaço quanto o grupo considerado nas ações educacionais, propondo campanhas para disseminar conhecimento sobre arte e/ou cultura para plateia maior do que a de estudantes.
De modo geral, a fluidez do texto das respostas e a consistência das ideias articuladas ficou aquém do padrão desejado para estudantes egressos do ensino superior.
O maior volume de respostas mostrou que os estudantes valorizaram mais a apresentação de duas ações educativas, muitas vezes de forma incompleta, apenas citando uma ação proposta e deixando de explicar o desdobramento solicitado no comando.
Além disso, houve respostas em que sugeriram, de forma isolada, ações do tipo apresentações, shows, exposições, sem qualquer detalhamento.
Entretanto, uma ação educativa com as finalidades solicitadas precisaria ser minimamente descrita para que fosse possível identificar que a ação pressupunha a presença de um mediador ( professor / arte educador / museólogo ) e de receptores participativos.
A proposição de uma ação educativa também foi confundida pelos estudantes com a proposta de implantação de equipamentos de infraestruturas ou projeto governamental, tais como : escolas, centros educacionais, incentivo ao ensino das artes em áreas mais necessitadas e introdução de programas de ensino de artes nos currículos escolares.
O contexto da arte e/ou da cultura e seus confrontos com a censura foram considerados, muitas vezes, sem o estabelecimento de relações, mas restringindo -se a aspectos descritivos de como o respondente compreendia cada um desses elementos.
Há de se considerar, adicionalmente, que em uma parte efetiva das respostas os estudantes defendiam, indevidamente, a censura como um mecanismo legitimamente regulador dos chamados"excessos da arte", revelando confusão entre os limites legais estabelecidos em lei ( que tipificam discursos de ódio, pedofilia, apologia às drogas, etc. ) e os julgamentos pessoais, coletivos ou institucionais que desencadeiam atitudes censoras.
A abrangência proporcionada pelo tema e pelo comando refletiu -se nas respostas, principalmente nos exemplos das ações educativas apresentados pelos estudantes.
Já na reflexão sobre arte, cultura e censura não houve tanta diversidade – a discussão quase sempre permanecia no senso comum, com menção à influência da época militar ou associada à discussão atual sobre limitações à liberdade de expressão.
Houve grande incidência de respostas em que o elemento censura foi confundido com a noção de classificação indicativa.
O desempenho ficou aquém do esperado.
Na maioria das respostas, os alunos ficaram com notas baixas, e foi grande a incidência de zeros.
Cabe ressaltar também a grande quantidade de respostas em branco e as que receberam nota zero devido a anulações por defesa de atos de censura, na contramão do que estabelecem a Constituição Federal e os Direitos Humanos.
A fuga ao que foi solicitado pelo comando da questão foi marcante, talvez evidenciando dificuldade de interpretação ou distanciamento em relação ao tema.
Outra lacuna muito observada foi a dificuldade de se expressar em textos fluentes e articulados.
A fragmentação e a falta de profundidade da argumentação foram uma constante.
Além disso, chamou a atenção a predominância do lugar comum na construção das respostas.
Tratando -se de um universo de formandos de nível superior, a superficialidade da reflexão em torno da arte e da censura e a dificuldade para incluir a Constituição Federal na discussão foram surpreendentes. 6.3.1.3 ANÁLISE DE CONTEÚDO DA QUESTÃO DISCURSIVA 2 DO COMPONENTE DE FORMAÇÃO GERAL A Tabela 6.15, mostra que o desempenho médio dos estudantes de Licenciatura, na Questão Discursiva 2 ( Média 29,1 ), foi maior ao obtido na Questão Discursiva 1 ( Média 22,5 ).
A região Norte foi aquela cuja Média foi menor ( 18,3 ), e a maior Média foi na região Sudeste ( 32,3 ).
Quanto à variabilidade das notas, o Desvio padrão de todo o Brasil, entre os cursos de Licenciatura, foi 25,5, superior ao obtido na Questão Discursiva 1 ( 19,0 ).
O maior desvio nessa questão foi obtido na região Sudeste ( 27,1 ), enquanto o menor foi obtido na região Sul ( 21,8 ).
A Mediana de todo o Brasil, entre os cursos de Licenciatura, foi 25,0, a mesma de três regiões.
As exceções foram a região Norte, que obteve Mediana 10,0 e a Centro-Oeste que obteve Mediana 30,0.
As notas : Máxima ( 100,0 ) e Mínima ( 0,0 ), foram as mesmas nas regiões Nordestes e Sudeste.
Já as demais regiões obtiveram nota Máxima e Mínima iguais a 95,0 e 0,0.
Entre os alunos de Bacharelado, o desempenho médio dos estudantes, na Questão Discursiva 2 ( Média 21,9 ), foi superior ao obtido na Questão Discursiva 1 ( Média 17,1 ).
A região Norte foi aquela cuja Média, também nessa questão, foi menor ( 14,5 ), mas a de maior Média nessa questão foi a região Sul ( 25,1 ).
Quanto à variabilidade das notas, o Desvio padrão de todo o Brasil, entre os cursos de Bacharelado, foi 22,5, superior ao obtido na Questão Discursiva 1 ( 18,0 ).
O maior desvio nessa questão foi obtido na região Centro-Oeste ( 23,3 ), enquanto o menor foi obtido na região Norte ( 19,3 ).
A Mediana de todo o Brasil, entre os cursos de Bacharelado, foi 15,0, a mesma da região Sudeste.
Para as demais regiões, a Mediana foi : 0,0 para região Norte, 10,0 para a região Nordeste, 30,0 para a região Sul e 25,0 para a região Centro-Oeste.
A nota Máxima ( 95,0 ) foi obtida por três regiões : Sudeste, Sul e Centro-Oeste, enquanto a região Norte obteve 65,0 e Nordeste obteve 80,0.
A nota Mínima zero foi igual para as cinco regiões.
No Gráfico 6.8, mostra -se a distribuição das notas na Questão Discursiva 2 do Componente de Formação Geral.
Observa -se que 12,2 \% dos estudantes de Licenciatura deixaram a questão em branco.
Chega a 24,5 \%, se agregado à frequência daqueles que, tendo respondido, receberam nota zero.
Nota -se, ainda, que as notas ficaram mais dispersas em comparação às da Questão Discursiva de número 1, o que pode ser constatado, também, pela comparação do Desvio padrão das notas da Questão Discursiva 2 ( 25,5 ) e o das notas da Questão Discursiva 1 ( 19,0 ).
Para a Licenciatura, a moda da distribuição das notas da Questão Discursiva 1 pode ser observada no intervalo ( 20 ; 30 ], com frequência de 23,1 \%.
Dentre os alunos de Bacharelado, observa -se que 10,8 \% dos estudantes deixaram a resposta à questão em branco.
Chega a 36,7 \%, se agregado à frequência daqueles que, tendo respondido, receberam nota zero, a classe modal.
Um máximo local pode ser observado no intervalo ( 20 ; 30 ], com frequência de 22,9 \%.
Nota -se, ainda, que as notas ficaram mais dispersas em comparação às da Questão Discursiva de número 1, o que pode ser constatado, também, pela comparação do Desvio padrão das notas da Questão Discursiva 2 ( 22,5 ) e o das notas da Questão Discursiva 1 ( 18,0 ). 6.3.1.4 COMENTÁRIOS SOBRE A CORREÇÃO DE CONTEÚDO DAS RESPOSTAS À QUESTÃO DISCURSIVA 2 O enunciado da segunda Questão Discursiva de Formação Geral era composto por dois textos motivadores e um comando com dois itens.
O primeiro texto apresentava uma definição de cidade inteligente, pelo uso da tecnologia nos diversos setores da administração e seu compromisso com o meio ambiente.
O segundo texto ampliava a definição anterior explorando aspectos da gestão pública.
O enunciado era claro, empregava linguagem adequada e continha informações suficientes para que os egressos construíssem seus argumentos.
O comando solicitava, primeiramente, que o estudante explicasse de que modo as cidades inteligentes podem contribuir para a melhoria das questões relacionadas ao desenvolvimento sustentável.
Em atendimento ao segundo item do comando, o estudante deveria apresentar uma proposta de intervenção urbana que gerasse impacto social e contribuísse para a melhoria da vida em comunidade.
Os dois itens do comando exigiam articulação de ideias e exposições com desdobramentos do porquê da resposta, o que tornou a questão complexa.
A questão abordava, portanto, temas previstos nas Diretrizes Curriculares para a Formação Geral : ciência, tecnologia e inovação ; biodiversidade, sustentabilidade e intervenção humana e cidades, habitação e qualidade de vida.
Além disso, avaliava o desenvolvimento de competências também previstas nas Diretrizes Curriculares para a Formação Geral : planejar e elaborar projetos de ação e intervenção a partir da análise de necessidades em contextos diversos ; ler, produzir e interpretar textos com clareza e coerência ; identificar, compreender e analisar situações-problema utilizando pensamento lógico e sistêmico, estabelecendo relações de causalidade.
Os conhecimentos e o nível de profundidade exigidos foram considerados adequados para egressos de cursos superiores pelos componentes da banca de correção.
Os textos tratavam de temas atuais e de amplo debate na mídia, tais como : cidade inteligente, meio ambiente, desenvolvimento sustentável, intervenção urbana e impacto social.
No entanto, o nível de dificuldade foi considerado elevado, já que tais conceitos admitem interpretações diversas nesse universo particular de egressos, o que criou, em consequência, um espectro diversificado de respostas e tipos de argumentação.
Além disso, era necessário conhecer o conceito de intervenção urbana para propor uma ação deste tipo que gerasse impacto social.
O padrão de resposta estava redigido de forma clara e, para o item ‘ a ’ do comando, apresentava um exemplo simples e informativo de como a tecnologia pode ser usada para a redução do consumo de energia e emissão de CO2.
Para o item ‘ b ’, o padrão apresentava exemplos de propostas de intervenção de forma mais diversificada do que a resposta padrão do item ‘ a ’, levantando várias opções corretas.
A abordagem adotada pelos estudantes foi quase sempre a da exemplificação, o mesmo recurso adotado no padrão de resposta.
Assim, pode -se dizer que as respostas dos estudantes se mantiveram, quase sempre, em conformidade com o contexto do padrão, mas com uma variedade muito grande de temas e ambientes relacionados ao desenvolvimento sustentável.
Dentre as ações vinculadas às cidades inteligentes, as ocorrências mais significativas foram as relacionadas à poluição, ao desmatamento, aos transportes coletivos, à segurança e ao desenvolvimento de tecnologia de modo a contribuir para a melhoria das dificuldades enfrentadas atualmente nas cidades.
Em muitos casos, o estudante articulou a sua resposta com a realidade vivenciada em determinadas cidades, comparando, inclusive, cidades que são consideradas inteligentes com realidades locais.
Também foram encontradas respostas em que os estudantes apresentavam propostas e soluções para realidades específicas das cidades onde o respondente reside, contextualizando o uso da tecnologia com ações fundamentais que faltam em algumas cidades brasileiras, como saneamento básico, água não poluída para uso doméstico, destinação do lixo.
As respostas incompletas ocorreram, muitas vezes, em função de o aluno apenas apresentar ações ou características das cidades inteligentes sem articular com o desenvolvimento sustentável ou com a intervenção urbana.
Em muitas respostas, os alunos apenas repetiram ideias do enunciado ; em outras, demonstraram a não compreensão sobre o que estava sendo solicitado pelo comando, deixando de argumentar sobre uma proposta ou ação mencionada.
Nas respostas, os estudantes se concentraram no uso da tecnologia para monitorar o desmatamento e a violência ; em hortas comunitárias para melhorar a qualidade da alimentação das pessoas ; em investimento em locais públicos para a realização de exercícios ; em maior controle do uso de combustíveis não fósseis para evitar a poluição.
As propostas de intervenção fixaram -se bastante em aplicativos, solução sugerida pelo ambiente atual de desenvolvimento e das facilidades proporcionadas pela Internet.
Marcante foi a diversidade de funções encontradas nas propostas de sensores e aplicativos, desde marcador de horário de passagem de ônibus até controlador de desperdício de algum bem, como água, por exemplo.
O resultado ficou abaixo do esperado.
A quantidade expressiva de zeros talvez tenha sido consequência de desinteresse mais do que por dificuldades para interpretar e responder ao comando.
Essas ficaram mais evidentes nas propostas de intervenção urbana, muitas vezes incompletas por sequer explicitarem o benefício delas advindo.
As maiores fragilidades reveladas pelas respostas dos estudantes, apontadas pelos corretores, foram a interpretação do contexto e do comando da questão, além de uma grande dificuldade para se expressarem de forma fluente e concisa.
As respostas com geração de ideias próprias ocorreram em escala bem menor do que aquelas em que os estudantes usaram do recurso da paráfrase pobre sobre o texto do enunciado, revelando dificuldade para criar e articular ideias a respeito dos temas em foco.
A linguagem utilizada nas respostas ficou, muitas vezes, aquém do padrão desejado para estudantes egressos do ensino superior.
A clareza das ideias apresentadas e a articulação do conteúdo com o comando da questão raramente foram encontradas nas respostas no nível esperado.
A dificuldade mais evidente dos estudantes foi a de estabelecer relações e encadear ideias para completar um raciocínio iniciado ; foram muitas as respostas fragmentadas. 6.3.1.5 CONSIDERAÇÕES FINAIS SOBRE A CORREÇÃO DE CONTEÚDO DAS QUESTÕES DISCURSIVAS DO COMPONENTE DE FORMAÇÃO GERAL As questões discursivas da Área de Formação Geral abordaram temas atuais e proporcionaram a avaliação de competências importantes dentre as previstas pelas suas Diretrizes Curriculares.
A análise das respostas revelou, principalmente, dificuldade dos estudantes na interpretação dos enunciados – talvez pela linguagem dos textos motivadores ser mais complexa do que a do cotidiano – e na compreensão dos comandos das questões – que pode ser fruto de desatenção ou de desconhecimento mais profundo do contexto.
Esses fatos parecem ter sido os principais responsáveis pelo resultado aquém do esperado em termos de avaliação quantitativa.
Mesmo os estudantes que responderam em conformidade com o padrão demonstraram um universo de argumentação muito limitado, pobre em vocabulário, com o recurso frequente de procurar nos textos introdutórios os elementos para a construção de suas frases.
Isso criou uma peculiaridade interessante : grupos de respostas muito parecidas.
Encadear e relacionar argumentos também foi uma deficiência percebida.
As razões eram apresentadas de forma confusa, sem fluência e sem constituir unidade.
O conteúdo das respostas aponta uma preocupação com o viés social dos problemas abordados ; são bons exemplos as propostas bastante diversificadas de melhoria de vida para a parcela da população menos favorecida.
Por outro lado, traz indícios de deficiências consideráveis na separação conceitual de cultura e arte, na compreensão mais ampla do conceito de censura e na capacidade de separar ações educativas de culturais ; em algumas respostas, também se evidenciaram conceitos equivocados sobre tecnologia, com exemplos de concretização muito difícil física ou financeiramente.
Houve alguns casos em que os conceitos de sustentabilidade e de meio ambiente se confundiam em lugar de se complementarem.
O \textbf{recorte} do grupo de formandos avaliado, em resumo, é de dificuldade de interpretação e produção de textos, vocabulário restrito ao praticado cotidianamente, assunção de fraseologia de senso comum em lugar do conteúdo apreendido formalmente e falta de profundidade em matérias importantes para a formação do cidadão como cultura, arte, 6.3.1.6 ANÁLISE DE LÍNGUA PORTUGUESA DAS QUESTÕES DISCURSIVAS DO COMPONENTE DE FORMAÇÃO GERAL Os dados obtidos a partir das respostas à Questão Discursiva do Componente de Formação Geral de Filosofia desagregados por Grande Região, segundo Habilitação e Modalidade de Ensino, no que tange à Língua Portuguesa, encontram -se na Tabela 6.16 e no Gráfico 6.9.
Nos cursos de Licenciatura em Filosofia, a Média desse aspecto foi 51,9, superior à obtida pelos alunos de Bacharelado ( 49,4 ).
Quando se considera a Modalidade de Ensino para cada Habilitação, os alunos de cursos Presenciais de Licenciatura obtiveram Média 51,7.
A Média com respeito à Língua Portuguesa foi maior para os alunos de cursos a Distância ( 52,1 ).
Para os cursos de Bacharelado, os valores foram : 49,0 para os cursos Presenciais, e 51,8 para os cursos a Distância.
A Mediana das notas de Língua Portuguesa para os cursos de Licenciatura foi 57,5 para ambas as Modalidades.
Para os cursos de Bacharelado, a Mediana foi 52,5, para o total e para os cursos Presenciais, e maior para o Bacharelado a Distância, 55,0.
A nota Máxima não chegou a 100,0 para nenhuma combinação de Habilitação e Modalidade.
No Gráfico 6.9, mostra -se a distribuição das notas de Língua Portuguesa das Questões Discursivas do Componente de Formação Geral.
Observa -se que a maior frequência 26,6 \%, corresponde à dos estudantes de Licenciatura que obtiveram nota no intervalo ( 50 ; 60 ].
Destacam -se, também, os estudantes de Licenciatura que deixaram a questão em branco, representando, aproximadamente, 7,7 \% do total.
Dentre os alunos de Bacharelado, mostra -se a distribuição das notas de Língua Portuguesa do Componente de Formação Geral da questão.
Observa -se que a maior frequência ( 29,4 \% ) corresponde à dos estudantes de Bacharelado que obtiveram nota no intervalo ( 50 ; 60 ].
Destacam -se, também, os estudantes de Bacharelado que deixaram a Discursiva 1 em branco, representando 7,4 \% do total. 6.3.1.7 COMENTÁRIOS SOBRE A CORREÇÃO DE LÍNGUA PORTUGUESA DAS RESPOSTAS ÀS QUESTÕES DISCURSIVAS A avaliação de desempenho linguístico De acordo com a portaria nº 386 do MEC/INEP, de 23 de agosto de 2021, o Enade deve avaliar se o estudante desenvolveu, no processo de formação, competências para : compreender as linguagens e respectivas variações ( item V ) ; ler, produzir e interpretar textos com clareza e coerência ( item VI ) ; analisar e interpretar representações verbais, não verbais de fenômenos diversos ( item VII ) ; identificar diferentes representações de um mesmo significado ( item VIII ) ; formular e articular argumentos e contra-argumentos consistentes em situações sócio comunicativas ( item IX ).
Essas competências são essenciais à formação de nível superior, porque dizem respeito ao domínio da linguagem, tanto no processo de leitura quanto na articulação e na organização de ideias para a produção textual.
Dessa forma, as duas questões discursivas de formação geral oferecem oportunidade para avaliar se o participante consegue desenvolver um tema, recorrendo a argumentos que explicitem sua opinião e elaborando propostas de intervenção social adequadas a um exercício consciente da cidadania, conforme determina a referida portaria : “ As questões discursivas do componente de Formação Geral avaliam aspectos como clareza, coerência, coesão, estratégias argumentativas, vocabulário e gramática adequados à norma-padrão da Língua Portuguesa. ” ( Art.2 §2 da referida portaria nº 386 ) Para avaliar esses aspectos, portanto, é necessário identificar, nas respostas discursivas de Formação Geral, as diferentes habilidades necessárias ao bom desempenho linguístico na situação de formalidade exigida : a organização das ideias, tendo em vista a construção de um texto coerente e estruturado ; o domínio das convenções ortográficas e das convenções morfossintáticas que constituem o arcabouço da gramática da norma-padrão da Língua Portuguesa ; e o domínio de um vocabulário amplo e diversificado, que permita uma comunicação clara e eficiente.
A Questão Discursiva 1 desenvolveu o tema da relação entre arte, cultura e censura, tomando como base dois fragmentos de textos.
O enunciado estava claro ao encaminhar a solicitação de resposta, evidenciando a exigência de elaboração de um texto dissertativo, por meio do verbo"discorrer", e de apresentação de duas ações educativas, articuladas com o ponto de vista desenvolvido na parte argumentativa.
A contraposição dos dois textos favoreceu a formulação e a articulação de argumentos e contra-argumentos consistentes, em defesa do ponto de vista a ser defendido pelo participante.
O enunciado permitia, inclusive, a ampliação da abordagem com base na análise de situações concretas que evidenciassem essa articulação.
A Questão Discursiva 2 desenvolveu o conceito de"cidade inteligente", tomando como base dois fragmentos de estudos que descreviam as suas características.
Com base nesses dois fragmentos, o comando solicitou do participante :"explique a contribuição das cidades inteligentes para o desenvolvimento sustentável"e apresente uma proposta de intervenção urbana que contribua para a melhoria da vida em comunidade.
O comando “ explicar ” demandava a elaboração de um texto de base expositiva, sem desenvolvimento de argumentação favorável ou contrária à configuração de uma cidade inteligente e de seus benefícios para a sociedade.
Em relação à segunda parte, o comando exigia a elaboração de uma proposta de intervenção urbana, que deveria conter informações sobre os agentes da intervenção, os recursos aplicados e o detalhamento dessa ação.
A segunda questão, portanto, se diferenciava da primeira, por solicitarem estruturas textuais distintas : a Questão 1 se consolidava com base em um modo de organização argumentativo e a Questão 2 se consolidava com base em um modo de organização expositivo.
Esses dois estilos textuais ofereceram oportunidade de avaliar se os participantes estavam habilitados a desenvolverem textos com esses dois distintos objetivos comunicativos, o que seria esperado ao final do curso de graduação e necessário à sua vida profissional.
A análise do desempenho linguístico esperado de um estudante que está concluindo seu curso superior apoiou -se no princípio de que ele deveria utilizar conhecimentos sobre os temas e estruturar seus textos de acordo com as características do registro formal adequado à situação comunicativa – avaliação de conhecimentos ao final de um curso superior.
Essa configuração determina exigências quanto à adequação da seleção vocabular, ao desenvolvimento do conteúdo, à estruturação sintática dos períodos, à organização lógica das ideias, à utilização de procedimentos de encadeamento textual e de referenciação, à obediência às exigências morfossintáticas próprias da modalidade escrita da norma-padrão, ao respeito às regras ortográficas de acentuação gráfica e de grafia das palavras.
O padrão de resposta utilizado na avaliação A partir desses preceitos, foram considerados, no padrão de resposta, as seguintes competências : ( 1 ) estruturação textual ; ( 2 ) respeito às convenções ortográficas da norma-padrão da Língua Portuguesa ; ( 3 ) domínio dos diferentes aspectos morfossintáticos próprios da modalidade escrita formal da norma-padrão da Língua Portuguesa ; e ( 4 ) seleção vocabular adequada à modalidade escrita formal da Língua Portuguesa, exigida pela situação comunicativa.
A seguir, os aspectos avaliados por cada uma destas competências, são apresentados. ( 1 ) Estruturação textual condizente com o gênero textual solicitado e o modo de organização textual adequado ao gênero ( argumentativo ou expositivo ).
Esta competência envolve : - estruturação sintática condizente com o padrão da modalidade escrita formal da Língua Portuguesa de modo a garantir a clareza e a coerência necessárias ; - distribuição do conteúdo do texto em parágrafos, de modo a garantir a sua organização temática ; - distribuição das ideias em períodos que mantenham uma estruturação sintática condizente com a modalidade formal da língua ( evitando -se frases fragmentadas, sequência justaposta de ideias sem encaixamentos sintáticos e presença de oração subordinada sem oração principal que reflitam hábitos da oralidade e comprometam a estrutura lógico-gramatical do texto ) ; - utilização de operadores discursivos que contribuam para a progressão temática do texto, estabelecendo relações lógicas entre as ideias apresentadas, tanto do ponto de vista intrafrasal, como interfrasal : preposições, conjunções, pronomes relativos, advérbios e locuções adverbiais de tempo e lugar ; - utilização de procedimentos de referenciação lexical e pronominal que permitam a retomada de referentes textuais, como pronomes, advérbios e artigos ; além de elementos lexicais, como sinônimos, hiperônimos, hipônimos, expressões metafóricas e expressões resumitivas ) ; - emprego adequado dos pronomes relativos “ cujo(a) ” e “ onde ”, responsáveis por encaixamentos sintáticos no período composto ; - utilização adequada de sinais de pontuação que contribuam para a organização lógica da frase e do texto ; - inteligibilidade relacionada ao atendimento das exigências de estruturação textual.
A expectativa, portanto, era de que o participante recorresse a procedimentos linguístico-discursivos para organizar seu texto, permitindo o encadeamento lógico entre suas partes de forma a garantir a progressão e a coerência textuais, com estruturação que não reproduzisse hábitos da oralidade e sim atendesse à norma-padrão da Língua Portuguesa. ( 2 ) Respeito às convenções ortográficas da norma-padrão da Língua Portuguesa.
Esta competência envolve o respeito às convenções estabelecidas pela legislação em vigor e consubstanciadas no Vocabulário Ortográfico da Língua Portuguesa, editado pela Academia Brasileira de Letras e avaliam -se : - domínio das regras de acentuação gráfica ; - domínio da grafia padrão das palavras ( com ausência de abreviaturas próprias da linguagem da internet ).
Assim, esperava -se que o participante : grafasse corretamente as palavras ; respeitasse as regras de acentuação gráfica ; empregasse maiúsculas em início de frase, em nomes próprios de pessoas, lugares ou instituições ; grafasse as siglas com letras maiúsculas ou apenas com a primeira letra maiúscula quando formarem uma palavra ; evitasse abreviações como p/, vc, tb, pq, tá, né, usadas muitas vezes em escrita informal e na internet ; obedecesse às regras de separação de sílabas no final da linha ; atendesse aos padrões de segmentação vocabular da norma-padrão. ( 3 ) Domínio dos diferentes aspectos morfossintáticos próprios da modalidade escrita formal da norma-padrão da Língua Portuguesa.
Esta competência envolve o domínio das convenções próprias da norma padrão da Língua Portuguesa relativas aos seguintes aspectos : concordância nominal ; concordância verbal ; regência nominal ; regência verbal ; flexão nominal ; flexão verbal ; correlação entre os tempos verbais ; colocação pronominal.
Assim, esperava -se que o participante : - flexionasse o verbo para estabelecer concordância de número com o sujeito da frase ; - flexionasse o artigo, o adjetivo e o pronome para concordar em número e em gênero com o substantivo a que cada um se refere ; - observasse a regência nominal e a verbal, utilizando a preposição adequada depois de um substantivo, um verbo ou um adjetivo ; - empregasse adequadamente o acento grave indicador da crase que expresse a contração entre uma preposição e um artigo, de acordo com as exigências de regência nominal ou verbal ; - obedecesse às regras de colocação pronominal ( próclise e ênclise ), distintas dos hábitos da oralidade ou da escrita informal ; - flexionasse adequadamente verbos, substantivos, adjetivos e pronomes, no que diz respeito à expressão das categorias gramaticais ; - flexionasse os verbos para expressar a correlação de modo e de tempo nas estruturas subordinadas.
Com base nesses critérios, alguns desvios ortográficos foram considerados como desvios de caráter morfossintático porque interferem na classificação dos vocábulos de base morfológica e funcional : - eliminação da marca de infinitivo ( -r- ) e substituição por acento agudo ou ausência total de marca do infinitivo, causando ambiguidade entre distintas classes de palavras ; - confusão entre as terminações “ ão ” e “ am ”, utilizadas para marcar diferentes tempos verbais ; - emprego indevido da forma verbal de"haver"\&lt ; “ há”\&gt ; no lugar da preposição ou do artigo \&lt ; “ a”\&gt ;, e vice-versa ; - emprego de hífen para separar pronome átono – tanto uso indevido quanto omissão ( ex : “ esperasse ”, em lugar de “ espera -se ” ; “ falar -mos ”, no lugar de “ falarmos ” ) ; - emprego de “ esta ” no lugar de “ está ”, provocando ambiguidade devido à confluência gráfica entre o pronome demonstrativo e a forma do verbo"estar", pela ausência da acentuação gráfica ; - emprego da conjunção aditiva"e"no lugar da forma verbal “ é ”, provocando ambiguidade devido à confluência gráfica por ausência da acentuação gráfica ; - emprego do advérbio “ mais ” no lugar da conjunção adversativa “ mas ”, reproduzindo confluência fônica da oralidade ; - registro da terceira pessoa do plural dos verbos “ ter ” e “ vir ” sem o acento circunflexo, causando desvio de concordância verbal ; - emprego inadequado do relativo “ onde ” deslocado do seu sentido locativo, reproduzindo tendência da oralidade ; - emprego do pronome relativo “ que ” sem a necessária preposição ( desvio de regência ) ; - ausência do acento grave obrigatório para marcar a crase. ( 4 ) Seleção vocabular adequada à modalidade escrita formal da Língua Portuguesa, exigida pela situação comunicativa.
Esta competência envolve : - precisão na seleção/utilização do vocabulário relacionado à temática solicitada pela questão ; - diversificação e adequação vocabular ; - coerência na combinação vocabular ; - ausência de marcas da oralidade, como termos de sentido muito genérico ( “ coisa ”, “ negócio ”, “ você ” ) e termos de registros mais informais ( como gírias, jargões, frases feitas, ditados populares, termos regionais ).
Assim, esperava -se que o participante : respeitasse a adequação vocabular, não empregando gírias ou expressões coloquiais ; apresentasse repertório vocabular diversificado e preciso ; evitasse repetição desnecessária de palavras ; utilizasse um vocabulário mais formal, como solicitado por um texto dissertativo.
Essas quatro competências foram reunidas em três itens de avaliação, atendendo aos aspectos explicitados anteriormente : I - Aspectos Ortográficos ( 20 \% ) : domínio das convenções ortográficas : grafia de vogais e consoantes, uso de maiúsculas e minúsculas, emprego do hífen e da acentuação gráfica, segmentação de palavras ; II - Aspectos Textuais ( 40 \% ) : domínio dos procedimentos de estruturação textual : organização interna dos períodos, emprego de conectores para a articulação lógica entre os períodos e entre os parágrafos, emprego de marcas de referenciação lexical e pronominal ; utilização dos sinais de pontuação que contribuem para a organização lógica da frase ; III - Aspectos Morfossintáticos e Voc vocabulares ( 40 \% ) : domínio das regras de caráter morfossintático estabelecidas como modelares do ponto de vista da modalidade escrita formal da norma-padrão da Língua Portuguesa : concordância nominal e verbal, regência nominal e verbal, colocação pronominal, flexão nominal e verbal, correlação entre tempos e modos verbais, ausência de marcas de oralidade.
Outra competência combinada a esta na análise das respostas é a seleção vocabular adequada à modalidade escrita formal da Língua Portuguesa, tendo em vista a intersecção entre as duas do ponto de vista das exigências do registro formal da modalidade escrita da norma-padrão.
Análise das respostas dadas pelos estudantes A Questão 1 propiciou a oportunidade de maior desenvolvimento do tema solicitado, gerando consequências na estruturação textual das respostas, já que os textos foram mais opinativos, mais elaborados e fluentes.
Já a Questão 2 teve um comportamento distinto, devido ao seu caráter mais expositivo, os textos foram mais curtos e o vocabulário mais limitado.
Observaram -se, nas duas questões, alguns casos de comprometimento de autoria, com textos muito próximos a paráfrases dos textos motivadores e sem marcas de repertório sociocultural produtivo.
Apresentam -se, a seguir, algumas observações sobre o desempenho dos participantes : I - Aspectos Ortográficos : o desempenho dos participantes revelou uma diferença muito grande nos dois aspectos analisados, como já foi observado nos relatórios dos anos anteriores : baixo índice de desvios da grafia padrão e grande índice de desvios de acentuação.
Em vários casos, ocorre ausência quase completa de acentuação gráfica.
Vale ressaltar que muitos participantes apresentaram uma caligrafia quase ininteligível.
Ainda nesse aspecto, ocorreu a omissão inadequada de cedilha, de til e do traço do T em algumas respostas.
Houve, inclusive, respostas com ausência total de cedilha ou til.
Os casos mais sistemáticos de eliminação do acento indicador da sílaba tônica foram : palavras proparoxítonas ( “ artistica ”, “ ludico ”, “ publico ”,"musica") ; palavras paroxítonas terminadas em ditongo crescente ( “ necessario ”,"historia","propria", “ vivencias ” ) ; palavras oxítonas ( “ tambem ”,"populacao").
Por outro lado, destaca -se o uso indevido do acento gráfico em palavras como “ propíciar ” “ criânça ”, por exemplo.
Quanto aos desvios de grafia, observam -se os seguintes exemplos : emprego inadequado de maiúscula : “ a Liberdade artística ” ; ou de minúscula :"constituição"; confusão entre"ss","sc","c"e"ç": “ condenassão ”, “ nessessario", “ repreiensão ”, “ conciencia ” ; eliminação de grafemas : “ essescia ”,"questõe"; entre"ss","sc","c"e"ç": “ condenassão ”, “ nessessario", “ repreiensão ”, “ conciencia ” ; eliminação de grafemas : “ essescia ”,"questõe"; - desvios de segmentação : “ ser vidor", “ igual dade ”, “ possi bilidades"; troca de grafema : “ convença ” ( por “ conversa ” ), “ comun ”.
II - Aspectos Textuais : a estruturação textual foi a competência que apresentou o maior número de desvios, sendo responsável, em alguns textos, pela falta de clareza e pela desorganização.
No conjunto dos textos analisados, observaram -se diferentes níveis de domínio dos recursos de estruturação e organização textual.
No entanto, os participantes podem ser agrupados em três faixas descritas a seguir : - os que possuem grande domínio dos recursos de estruturação textual, com produção de textos que atendem às exigências de clareza e coerência, com progressão temática bem direcionada e utilização produtiva de operadores argumentativos ; - os que apresentam um domínio regular dos recursos de estruturação textual, com desvios na organização do texto e ausência de paragrafação, além de uso insuficiente de operadores argumentativos e processos de encaixamento sintático ; - os que apresentam baixo domínio dos recursos de estruturação textual, com muitos desvios que retomam hábitos da oralidade e que comprometem a coerência do texto : ocorrência de sequência justaposta de ideias sem encaixamentos sintáticos ; redução de estruturas subordinadas, ao lado do aumento na frequência de estruturas coordenadas e absolutas ; redução no uso de conectores para expressar relações lógicas essenciais à construção do texto ; emprego equivocado de operadores que não estabelecem relações lógicas coerentes entre ideias do texto ; fragmentação frasal que compromete a estrutura lógico-gramatical.
Uma estrutura que vem se ampliando na escrita dos participantes, refletindo hábitos da oralidade, é a duplicação de um dos elementos da frase, ou seja, a retomada anafórica de um elemento topicalizado, como por exemplo :"Mesmo com a época da censura, o ser humano ELE pode se expressar de várias formas..."; “ A cultura ELA vem trazida de várias épocas... ” ; “ Arte e cultura, ELAS estão relacionadas em uma só ” ; “ A cultura ELA vem enraizada a pessoa... ” ; “ As ações educativas ELAS podem ser consideradas boas ou ruins. ” Também em relação ao emprego dos sinais de pontuação, observaram -se distintos tipos de desempenho : participantes que dominam a funcionalidade da pontuação para a construção de textos fluentes e coesos, por um lado, e um grande número de participantes que apresentam precariedade em sua utilização, chegando ao ponto de elaborar parágrafos sem qualquer marca interna de pontuação para separar os períodos.
São os seguintes os tipos de desvios encontrados : - emprego de vírgula para separar o sujeito e o predicado : “ Os teatros, festas de rua, carnaval, elevam nosso conhecimento ” ( no lugar de “ Os teatros, festas de rua, carnaval elevam nosso conhecimento ” ) ;"que governos e municípios, façam a inclusão ” ( no lugar de “ que governos e municipios façam a inclusão ” ) ; - ocorrência de apenas uma das vírgulas para separar uma palavra, uma expressão ou uma oração encaixada : “ A liberdade artística que está presente em nossa Constituição 88, nos calça ” ( no lugar de “ A liberdade artística, que está presente em nossa Constituição 88, nos calça ) ” ; - emprego de ponto no lugar de vírgula para separar orações de um mesmo período : “ criando debate em espaços públicos.
Além de convidar os críticos a mostrarem o que é arte! ” ( no lugar de “ criando debate em espaços públicos, além de convidar os críticos a mostrarem o que é arte! ” ) ; - ausência de vírgula ou conjunção para separar elementos de uma enumeração :"As cidades inteligentes são caracterizadas por inteligência artificial, desenvolvimento sustentável conservação dos espaços verdes"( no lugar de"As cidades inteligentes são caracterizadas por inteligência artificial, desenvolvimento sustentável e conservação dos espaços verdes") ; - ausência de vírgula para separar oração adjetiva explicativa ou utilização inadequada para separar oração adjetiva restritiva :"As cidades inteligentes que são lugares em que a tecnologia avançada é utilizada..."( no lugar de"As cidades inteligentes, que são lugares em que a tecnologia avançada é utilizada,...") ; - emprego indevido do sinal de dois pontos : “ percebe -se um círculo onde estão infinitamente ligadas : a cultura e a arte ” ( no lugar de “ percebe -se um círculo onde estão infinitamente ligadas a cultura e a arte ” ) ; - ausência de pronome no processo de referenciação : “ censura a arte ou até mesmo impulsiona ” ( no lugar de “ censura a arte ou até mesmo a impulsiona ” ) ; - emprego de vírgula em lugar do ponto-final : “ são muitos os agravamentos em relação à tal situação, ” III - Aspectos Morfossintáticos e Vocabulares : o domínio das convenções morfossintáticas apresentou resultados que evidenciam, claramente, alguns aspectos problemáticos no desempenho dos participantes.
Assim como na análise dos desvios textuais, pode -se observar uma distinção entre três grupos de participantes : o grupo que revela domínio das convenções morfossintáticas ; o grupo que apresenta um domínio regular, com desvios pontuais ; e o grupo que revela domínio precário dessas convenções.
Pode -se destacar a regência, tanto verbal quanto nominal, como o aspecto com maior número de desvios : a ) Regência nominal e verbal : Na maioria dos textos, destacam -se os desvios relacionados à noção de"regência", ou seja, à relação de subordinação que se estabelece entre as palavras por meio, ou não, de preposições, refletindo uma tendência na modalidade escrita da língua.
O desvio mais frequente foi a falta do sinal indicativo da crase, em contextos em que seria obrigatória, ou, ao contrário, um uso indevido em contextos em que o princípio não se aplica.
Os dois tipos de desvios revelam que o usuário não tem consciência de que, sob a forma do termo “ a ”, se pressupõe a presença de uma preposição “ a ”, exigida pela regência do termo anterior.
Exemplos : “ que se contrapõem a essência das artes ” ( no lugar de “ que se contrapõem à essência das artes. ” ) ; “ Esse ato pode levar a naturalização"( no lugar de “ Esse ato pode levar à naturalização") ; “ Além de convidar os críticos a mostrarem o que é arte à fazendo! ” ( no lugar de “ Além de convidar os críticos a mostrarem o que é arte a fazendo! ” ).
Outro desvio relacionado à regência verbal e à nominal, encontrado frequentemente nas questões, é a ausência de preposição antes de pronome relativo, refletindo um processo generalizado na modalidade oral da língua, em situações de registro informal.
Exemplo :"a cidade que todos querem morar é..."( no lugar de"a cidade em que todos querem morar é...") b ) Concordância nominal e verbal : Os desvios relativos à aplicação das regras de concordância verbal e nominal correspondem a fatos que têm sido amplamente estudados na modalidade oral da Língua Portuguesa.
Quanto à concordância de número, observou -se ausência de marca ( com sujeito anteposto ou posposto ), como por exemplo a ausência de acento circunflexo na forma plural do presente do indicativo do verbo “ ter ”, considerada como um desvio na concordância verbal e não na acentuação gráfica.
Exemplo :"os governantes tem procurado planejar"( no lugar de"os governantes têm procurado planejar") ; “ onde é mostrado os acontecimentos da atualidade ” ( no lugar de “ onde são mostrados os acontecimentos da atualidade ” ). c ) Registro da forma verbal “ é ” sem acento, provocando ambiguidade na leitura pela confluência com a conjunção aditiva"e".
Exemplo : “ A principal forma de realizar essa disseminação e criar nas escolas... ” ( no lugar de “ A principal forma de realizar essa disseminação é criar nas escolas... ” ). d ) Emprego do pronome relativo “ onde ” como relativo universal, típica marca da oralidade, em que esse pronome não é mais utilizado apenas em sua função locativa"lugar onde se realiza a ação".
Exemplo :"a aceitação não é mais a mesma, onde muitas vezes é a cultura que censura a arte"; “ situações onde ” ; “ cultura mais cotidiana, onde ”. e ) Falta de artigo definido antes de substantivo, comprometendo o paralelismo em enumerações : Exemplo :"censurar a arte e cultura é prejudicial"( no lugar de"censurar a arte e a cultura é prejudicial"). f ) Desvios relacionados à flexão verbal : Exemplos : emprego indevido do tempo ou do modo verbal : “ adaptado ” por “ adaptando ” ; eliminação de marca de infinitivo : “ ta ” por “ estar ”, “ está ” por “ estar ”, “ esta ” por “ estar ” ; “ estás ” por “ estas ” ; utilização de construção de “ estar ” + gerúndio como marca de Outro desvio na aplicação da regra da concordância de número foi o aparecimento da marca de plural em verbos ou adjetivos relacionados a núcleos substantivos no singular, devido à presença de modificadores no plural, evidenciando um processo de hipercorreção.
Esse processo tem se ampliado nas situações de oralidade, caracterizando uma tendência de variação da Língua Portuguesa.
Exemplos :"a construção de teatros e cinemas revelam uma preocupação das cidades inteligentes"( no lugar de"a construção de teatros e cinemas revela uma preocupação das cidades inteligentes") ;"o aumento dos acidentes ambientais devem ser evitados ” ( no lugar de"o aumento dos acidentes ambientais deve ser evitado").
Quanto à concordância de gênero, vários casos foram observados, normalmente no âmbito de sintagmas nominais longos, em que o adjetivo está afastado do substantivo, conforme apontam as pesquisas sobre a modalidade oral.
Exemplos : “ sendo valorizado uma normatização de discurso que se contrapõe ” ( no lugar de “ sendo valorizada uma normatização de discurso que se contrapõe ” ) ; “ podendo essa agradar um público e desagradar outra ” ( no lugar de “ podendo essa agradar um público e desagradar outro ” ). c ) Registro da forma verbal “ é ” sem acento, provocando ambiguidade na leitura pela confluência com a conjunção aditiva"e".
Exemplo : “ A principal forma de realizar essa disseminação e criar nas escolas... ” ( no lugar de “ A principal forma de realizar essa disseminação é criar nas escolas... ” ). d ) Emprego do pronome relativo “ onde ” como relativo universal, típica marca da oralidade, em que esse pronome não é mais utilizado apenas em sua função locativa"lugar onde se realiza a ação".
Exemplo :"a aceitação não é mais a mesma, onde muitas vezes é a cultura que censura a arte"; “ situações onde ” ; “ cultura mais cotidiana, onde ”. e ) Falta de artigo definido antes de substantivo, comprometendo o paralelismo em enumerações : Exemplo :"censurar a arte e cultura é prejudicial"( no lugar de"censurar a arte e a cultura é prejudicial"). f ) Desvios relacionados à flexão verbal : Exemplos : emprego indevido do tempo ou do modo verbal : “ adaptado ” por “ adaptando ” ; eliminação de marca de infinitivo : “ ta ” por “ estar ”, “ está ” por “ estar ”, “ esta ” por “ estar ” ; “ estás ” por “ estas ” ; utilização de construção de “ estar ” + gerúndio como marca de Assim como nas demais competências avaliadas, há diferenças entre os participantes em relação à seleção e ao emprego do vocabulário.
Foram encontrados textos com alto padrão de domínio vocabular, chegando até ao emprego de vocabulário técnico das áreas do saber abordadas nas questões : arte, cultura, meio ambiente e tecnologia.
Por outro lado, na maioria dos textos foram observados vários tipos de inadequação na seleção e na utilização do vocabulário, sendo que o principal deles foi a excessiva repetição de certas palavras, revelando limitação de repertório vocabular.
Além disso, observou -se a ocorrência de expressões da oralidade, revelando falta de domínio de vocabulário mais abstrato e de maior complexidade, essencial ao desenvolvimento do texto de base dissertativa.
Exemplos : seleção vocabular indevida : “ \textbf{crescemos} tão intrínsecos a esses usos... ” ; emprego inadequado do verbo “ ter ” no lugar de “ haver ” : “ Durante a ditadura teve muita repressão ” ; imprecisão vocabular : “ ela anda lado a lado com a existência do que conhecemos de cultura ”, “ A liberdade artística que está presente em nossa Constituição 88, nos calça e... ”.
A imprecisão vocabular, como o uso da palavra “ coisa ” em vários textos, demonstra o domínio precário da norma-padrão da Língua Portuguesa, com presença constante de traços de oralidade. 6.3.1.8 CONSIDERAÇÕES FINAIS SOBRE A CORREÇÃO DE LÍNGUA PORTUGUESA DAS RESPOSTAS ÀS QUESTÕES DISCURSIVAS DO COMPONENTE DE FORMAÇÃO GERAL O conjunto de textos analisados no Enade 2021 permite traçar um perfil bem detalhado do desempenho linguístico do estudante de graduação brasileiro.
Conforme ficou comprovado nos comentários dos exemplos de respostas apresentados neste relatório, uma parcela significativa dos participantes revela lacunas importantes em sua formação, principalmente no que diz respeito à estruturação e à organização de textos.
Como esta é uma competência essencial à produção escrita, que deve ser construída ao longo dos anos de escolaridade ( do ensino fundamental ao ensino superior ), sua recuperação é trabalhosa e depende de um esforço intenso de exposição a textos e de reflexão sobre os requisitos que garantem a clareza e o encadeamento das ideias.
A competência relacionada ao domínio do vocabulário, que deve ser diversificado e apropriado ao padrão formal da língua, também apresentou um desempenho insatisfatório 6.3.2 COMPONENTE DE CONHECIMENTO ESPECÍFICO Na parte da prova relativa às questões discursivas do Componente de Conhecimento Específico ( Tabela 6.17a ), observa -se que a Média foi bem mais baixa do que para as questões discursivas do Componente de Formação Geral.
Enquanto, no Componente de Formação Geral, a Média para estudantes de Filosofia ( Licenciatura ) de todo o Brasil foi 30,9, na parte de Conhecimento Específico, a Média foi 30,7.
A maior Média deste componente foi obtida pelos estudantes da região Sul ( 32,4 ), e a menor, pelos da região Norte ( 27,2 ).
As demais regiões obtiveram Média igual a 27,8 ( Nordeste ), 31,9 ( Sudeste ) e 28,0 ( Centro-Oeste ).
Quanto à variabilidade das notas, o Desvio padrão de todo o Brasil, entre os alunos de Licenciatura, foi 25,1.
O maior Desvio padrão foi encontrado na região Centro-Oeste ( 26,8 ), e o menor, na região Norte ( 21,6 ).
A maior nota Máxima, 100,0, foi em todas as regiões com exceção da região Norte que obteve nota 85,0.
A nota Mínima zero foi obtida por, pelo menos, um estudante de Licenciatura pela maioria dos participantes, que se limitam ao senso comum, com vocabulário genérico e vago, o que é inadequado para atender ao perfil de um graduando.
As outras duas competências avaliadas ( convenções ortográficas, convenções morfossintáticas ) apresentam desvios pontuais e são de mais fácil recuperação, com exceção dos desvios morfossintáticos que já estão incorporados na modalidade oral informal da Língua Portuguesa, configurando processos de variação linguística e, talvez, antecipando futuras mudanças no padrão da modalidade escrita.
Com base na avaliação do desempenho linguístico aqui relatada, conclui -se que é preciso aperfeiçoar a formação superior, dando aos estudantes mais acesso às ferramentas necessárias para que sejam bem-sucedidos nas diferentes situações de sua vida cidadã – profissional, cultural e política.
Essa é uma contribuição efetiva que os resultados do Enade podem oferecer às instituições de ensino superior do país.
Para finalizar, vale observar que as duas questões discursivas atenderam às exigências de avaliação, tanto do domínio discursivo dos estudantes, quanto do seu conhecimento sobre temas relacionados ao conhecimento acadêmico, o que pode ser extremamente útil para que os cursos de ensino superior possam contribuir na formação de seus graduandos.
Se o objetivo é avaliar a competência discursiva dos formandos brasileiros, a elaboração de uma resposta de base dissertativo-argumentativa ( Questão 1 ) e de uma resposta de base expositiva ( Questão 2 ) certamente ofereceu material relevante para essa análise, por incentivar a reflexão e a organização do pensamento por meio de uma sequência textual coerente e complexa. 6.3.2 COMPONENTE DE CONHECIMENTO ESPECÍFICO Na parte da prova relativa às questões discursivas do Componente de Conhecimento Específico ( Tabela 6.17a ), observa -se que a Média foi bem mais baixa do que para as questões discursivas do Componente de Formação Geral.
Enquanto, no Componente de Formação Geral, a Média para estudantes de Filosofia ( Licenciatura ) de todo o Brasil foi 30,9, na parte de Conhecimento Específico, a Média foi 30,7.
A maior Média deste componente foi obtida pelos estudantes da região Sul ( 32,4 ), e a menor, pelos da região Norte ( 27,2 ).
As demais regiões obtiveram Média igual a 27,8 ( Nordeste ), 31,9 ( Sudeste ) e 28,0 ( Centro-Oeste ).
Quanto à variabilidade das notas, o Desvio padrão de todo o Brasil, entre os alunos de Licenciatura, foi 25,1.
O maior Desvio padrão foi encontrado na região Centro-Oeste ( 26,8 ), e o menor, na região Norte ( 21,6 ).
A maior nota Máxima, 100,0, foi em todas as regiões com exceção da região Norte que obteve nota 85,0.
A nota Mínima zero foi obtida por, pelo menos, um estudante de Licenciatura em todas as regiões do Brasil.
A Mediana do Brasil e das regiões Sul e Sudeste foi igual a 30,0 ; e nas regiões Norte, Nordeste e Centro-Oeste foi 25,0.
Na parte da prova relativa às questões discursivas do Componente de Conhecimento Específico ( Tabela 6.17b ), observa -se que a Média foi bem mais alta do que para as questões discursivas do Componente de Formação Geral.
Enquanto, no Componente de Formação Geral, a Média para estudantes de Filosofia ( Bacharelado ) de todo o Brasil foi 25,4, na parte de Conhecimento Específico, a Média foi 32,3.
A maior Média deste componente foi obtida pelos estudantes da região Sul ( 37,3 ), e a menor, pelos da região Norte ( 12,9 ).
As demais regiões obtiveram Média igual a 27,1 ( Nordeste ), 34,0 ( Sudeste ) e 27,4 ( Centro-Oeste ).
Quanto à variabilidade das notas, o Desvio padrão de todo o Brasil, entre os alunos de Bacharelado foi 25,2.
O maior Desvio padrão foi encontrado na região Sudeste ( 25,5 ), e o menor, na região Norte ( 19,4 ).
A maior nota Máxima, 100,0, foi obtida nas regiões Sul, Sudeste e Centro-Oeste.
Nas demais regiões, as notas Máximas foram : 80,0 ( Norte, a mais baixa ) e 90,0 ( Nordeste ).
A nota Mínima zero foi obtida por, pelo menos, um estudante de Bacharelado em todas as regiões do Brasil.
A Mediana do Brasil foi igual a 30,0 ; na região Norte foi 0,0 ; nas regiões Nordeste e Centro-Oeste foram 25,0 ; na região Sudeste foi 35,0 ; na região Sul, foi 40,0 – a maior.
No Gráfico 6.10, representa -se a distribuição das notas nas questões discursivas no Componente de Conhecimento Específico.
A moda dessa distribuição ocorre no intervalo [ 0 ; 10 ], para os alunos de Licenciatura, com 32,1 \% do total de participantes.
Observa -se que nesse intervalo estão computadas as ocorrências de respostas em branco para as três questões discursivas de Conhecimentos Específicos.
Dentre os alunos de Bacharelado, a moda dessa distribuição também ocorre no intervalo [ 0 ; 10 ], com frequência igual a 28,7 \%.
Observa -se que nesse intervalo estão computadas as ocorrências de respostas em branco para as três questões discursivas de Conhecimentos Específicos.
A análise de cada uma destas questões será feita a seguir. 6.3.2.1 ANÁLISE DA QUESTÃO DISCURSIVA 3 DO COMPONENTE DE CONHECIMENTO ESPECÍFICO A Questão Discursiva 3 de Conhecimentos Específicos foi a mesma para as duas habilitações.
Os resultados para a Licenciatura constam da Tabela 6.18a e para o Bacharelado estão apresentados na Tabela 6.18b.
Cabe esclarecer que por terem respondido à mesma questão os resultados são comparáveis.
Na Questão 3, cujos resultados aferidos se encontram descritos na Tabela 6.18a, a Média dos estudantes de Licenciatura de todo o Brasil foi 44,0.
O desempenho dos estudantes de Licenciatura nessa questão foi muito superior ao da Questão 5.
A menor Média nessa questão foi obtida pelos estudantes de Licenciatura da região Centro-Oeste ( 37,1 ), enquanto a maior Média foi obtida na região Sudeste ( 47,9 ).
Quanto à variabilidade das notas, o Desvio padrão dos estudantes de Licenciatura de todo o Brasil foi 35,7.
O maior Desvio padrão foi obtido nas regiões Nordeste e Centro-Oeste ( 35,9 ), enquanto o menor foi obtido na região Norte ( 32,7 ).
A nota Máxima, 100,0 pontos, foi alcançada por, pelo menos, um estudante de Licenciatura em cada região.
A Mediana dos estudantes de Licenciatura de todo o Brasil foi Na Questão 3, cujos resultados aferidos se encontram descritos na Tabela 6.18b, a Média dos estudantes de Bacharelado de todo o Brasil foi 46,0.
O desempenho dos estudantes de Bacharelado nessa questão foi muito superior ao verificado na Questão 5.
A menor Média nessa questão foi obtida pelos estudantes da região Norte ( 14,7 ), enquanto a maior Média foi obtida na região Sul ( 53,0 ).
Quanto à variabilidade das notas, o Desvio padrão dos estudantes de Bacharelado de todo o Brasil foi 35,8.
O maior Desvio padrão foi obtido na região Centro-Oeste ( 36,8 ), enquanto o menor foi obtido na região Norte ( 22,9 ).
A nota Máxima, 100,0 pontos, foi alcançada por, pelo menos, um estudante de Bacharelado em cada região, com exceção da região Norte, onde a nota Máxima foi 80,0.
A 40,0, para maioria das regiões ; na região Sudeste foi 50,0.
A nota Mínima zero foi a mesma em todas as regiões do Brasil.
As Médias para os alunos de Licenciatura ( ver Tabela 6.18a ) foram, respectivamente para os de cursos Presenciais e a Distância, 48,0 e 39,5 ( 44,0 para o total ).
A nota Máxima ( 100,0 ) foi atingida por, pelo menos, um concluinte em todas as regiões para ambas as Modalidades.
A Mediana foi maior para os cursos Presenciais ( 50,0 ) em comparação com os cursos a Distância ( 40,0 ).
As Medianas dos cursos Presenciais foram : 40,0, para as regiões Norte, Nordeste e Centro-Oeste ; 50,0 para a região Sul ; e 70,0 para a região Sudeste.
Já as Medianas dos cursos a Distância foram : 10,0 para a região Nordeste, 40,0 para as regiões Sudeste e Sul, 50,0 para região Centro-Oeste.
Cumpre notar que somente na região Norte não houve participação de estudantes de cursos a Distância de Licenciatura.
A nota Mínima foi zero para todas as combinações de Modalidade de Ensino e região.
Na Questão 3, cujos resultados aferidos se encontram descritos na Tabela 6.18b, a Média dos estudantes de Bacharelado de todo o Brasil foi 46,0.
O desempenho dos estudantes de Bacharelado nessa questão foi muito superior ao verificado na Questão 5.
A menor Média nessa questão foi obtida pelos estudantes da região Norte ( 14,7 ), enquanto a maior Média foi obtida na região Sul ( 53,0 ).
Quanto à variabilidade das notas, o Desvio padrão dos estudantes de Bacharelado de todo o Brasil foi 35,8.
O maior Desvio padrão foi obtido na região Centro-Oeste ( 36,8 ), enquanto o menor foi obtido na região Norte ( 22,9 ).
A nota Máxima, 100,0 pontos, foi alcançada por, pelo menos, um estudante de Bacharelado em cada região, com exceção da região Norte, onde a nota Máxima foi 80,0.
A Mediana dos estudantes de Bacharelado de todo o Brasil, foi 50,0, a mesma para a região Sudeste.
A Mediana das demais regiões foram : 0,0 ( Norte ), 40,0 ( Nordeste e Centro-Oeste ), 60,0 ( Sul ).
As Médias para os alunos de Bacharelado ( ver Tabela 6.18b ) foram, respectivamente para os de cursos Presenciais e a Distância, 45,5 e 48,5 ( 46,0 para o total ).
A nota Máxima ( 100,0 ) foi atingida por, pelo menos, um concluinte das regiões em todas as regiões e em ambas Modalidades, com exceção da região Norte na Modalidade Presencial onde a nota Máxima foi 80,0.
A Mediana foi 50,0 para os cursos de ambas as Modalidades.
Para os cursos Presenciais, as Medianas foram : 0,0 ( Norte ), 40,0 ( Nordeste e Centro-Oeste ), 50,0 ( Sudeste ) e 60,0 ( Sul ), enquanto para os cursos a Distância, foram : 50,0 ( Sudeste ), 55,0 ( Sul ) e 45,0 ( Centro-Oeste ).
Cumpre notar que somente nas regiões Norte e Nordeste não houveram participação de estudantes de cursos a Distância de Bacharelado.
A nota Mínima foi zero para todas as combinações de Modalidade de Ensino e região.
No Gráfico 6.11, mostra -se a distribuição das notas na Questão Discursiva 3, do Componente de Conhecimento Específico da Área de Filosofia ( Licenciatura/Bacharelado ).
Sem levar em consideração que 14,1 \% dos estudantes de Licenciatura deixaram a resposta à questão em branco e que 10,6 receberam nota zero, a moda dessa distribuição seria o intervalo ( 70 ; 80 ], com 21,8 \%, do total de participantes de Licenciatura.
Dentre os alunos de Bacharelado, a moda desta distribuição ocorre no intervalo ( 70 ; 80 ], com frequência igual a 22,9 \%.
Considerando -se que 15,0 \% dos estudantes de Bacharelado deixaram a resposta à questão em branco e que 9,7 \% receberam nota zero.
Mediana dos estudantes de Bacharelado de todo o Brasil, foi 50,0, a mesma para a região Sudeste.
A Mediana das demais regiões foram : 0,0 ( Norte ), 40,0 ( Nordeste e Centro-Oeste ), 60,0 ( Sul ).
As Médias para os alunos de Bacharelado ( ver Tabela 6.18b ) foram, respectivamente para os de cursos Presenciais e a Distância, 45,5 e 48,5 ( 46,0 para o total ).
A nota Máxima ( 100,0 ) foi atingida por, pelo menos, um concluinte das regiões em todas as regiões e em ambas Modalidades, com exceção da região Norte na Modalidade Presencial onde a nota Máxima foi 80,0.
A Mediana foi 50,0 para os cursos de ambas as Modalidades.
Para os cursos Presenciais, as Medianas foram : 0,0 ( Norte ), 40,0 ( Nordeste e Centro-Oeste ), 50,0 ( Sudeste ) e 60,0 ( Sul ), enquanto para os cursos a Distância, foram : 50,0 ( Sudeste ), 55,0 ( Sul ) e 45,0 ( Centro-Oeste ).
Cumpre notar que somente nas regiões Norte e Nordeste não houveram participação de estudantes de cursos a Distância de Bacharelado.
A nota Mínima foi zero para todas as combinações de Modalidade de Ensino e região.
No Gráfico 6.11, mostra -se a distribuição das notas na Questão Discursiva 3, do Componente de Conhecimento Específico da Área de Filosofia ( Licenciatura/Bacharelado ).
Sem levar em consideração que 14,1 \% dos estudantes de Licenciatura deixaram a resposta à questão em branco e que 10,6 receberam nota zero, a moda dessa distribuição seria o intervalo ( 70 ; 80 ], com 21,8 \%, do total de participantes de Licenciatura.
Dentre os alunos de Bacharelado, a moda desta distribuição ocorre no intervalo ( 70 ; 80 ], com frequência igual a 22,9 \%.
Considerando -se que 15,0 \% dos estudantes de Bacharelado deixaram a resposta à questão em branco e que 9,7 \% receberam nota zero. 6.3.2.2 COMENTÁRIOS SOBRE AS RESPOSTAS À QUESTÃO DISCURSIVA 3 A Questão Discursiva de número 3 foi proposta para estudantes do Bacharelado e da Licenciatura em Filosofia.
O tema contemplado foi o conceito de indústria cultural que é central no pensamento de T.
Adorno e M.
Horkheimer, e uma das grandes contribuições teóricas dos filósofos frankffurtianos para a análise da sociedade contemporânea.
Os projetos \textbf{pedagógicos} dos cursos de Bacharelado e Licenciatura em Filosofia orientam -se pelas Diretrizes Curriculares da área para fornecerem formação sólida da história da Filosofia aos estudantes, a fim de que eles consigam, dentre outros aspectos, compreender e transmitir diferentes tipos de problemas e temas e também distintos sistemas filosóficos.
Para tanto, torna -se necessário o desenvolvimento de competências e habilidades relacionadas às seguintes capacidades : ler, interpretar, analisar e comentar textos teóricos ; compreender o sentido da existência e das produções culturais ; perceber a integração necessária entre a Filosofia e a produção artístico-cultural ; analisar criticamente a realidade sócio-histórico-política em que se insere.
O enunciado dessa questão apresentava um pequeno trecho do capítulo"A indústria cultural : o iluminismo como mistificação de massas", dos filósofos Adorno e Horkheimer, e o comando solicitava a resposta dividida em duas partes.
Foi solicitado que o/a estudante explicasse o conceito frankfurtiano de indústria cultural ( item ‘ a ’ ) e também analisasse como a indústria cultural molda a realidade dos/das espectadoras ( item ‘ b ’ ).
O enunciado da questão era claro e objetivo e também continha informações suficientes para a construção das respostas solicitadas.
Sendo assim, não havia possibilidade de interpretações diferentes do esperado pelo padrão.
Em geral, a diversidade de respostas equivocadas, erradas ou que fugiram ao padrão refletiam, basicamente, desconhecimento do tema.
O nível de conhecimento exigido pela questão era adequado aos formandos da área e condizente com o que costuma ser discutido em disciplinas como Estética, Filosofia da Arte, Filosofia Contemporânea, Filosofia Social, Filosofia Política, Filosofia da Cultura, uma vez que o solicitado era básico para o entendimento do conceito criado pelos filósofos da Escola de Frankfurt.
Esse cenário é ratificado pela importância que o conceito de indústria cultural possui na atualidade, notadamente marcada pela massiva circulação de produtos culturais por diferentes mídias, linguagens e formatos.
Nesse sentido, conhecer o conceito e ter a capacidade de operacionalizá -lo na análise de situações contemporâneas vai ao encontro das competências e habilidades requeridas do egresso, segundo as Diretrizes Curriculares para os Cursos de Graduação em Filosofia.
Por esses motivos, considerou -se a questão fácil e de baixa complexidade.
O padrão de resposta era claro e detalhado, refletindo a preocupação com a diversidade de argumentos que poderiam ser utilizados pelos estudantes, o que contribuiu para a correção da questão.
Via de regra, nas respostas corretas ( em parte ou integralmente ), os estudantes não fugiram das possibilidades presentes no padrão de resposta, em que se contemplou boa diversidade de caminhos argumentativos possíveis.
As respostas incorretas ou equivocadas ou deixadas em branco resultaram do desconhecimento sobre o tema"indústria cultural"e/ou do conceito e da abordagem frankfurtiana.
Foi percebida diferença qualitativa em relação à linguagem adotada por grupos de estudantes na redação da resposta, sem que fosse possível identificar tais grupos com apenas uma das habilitações : Bacharelado ou Licenciatura.
Excluindo -se as notas zero oriundas de respostas em branco, desconsideradas ou nulas, os percentuais de notas entre zero e cinco, em torno de 47 \% no Bacharelado, e de 52 \% na Licenciatura, estão associados ao fato de parte considerável dos/das respondentes utilizar linguagem bem primária e simplória, próxima do senso comum, o que demonstra baixa proficiência na exposição de ideias e na construção de argumentos.
Em geral, foi percebido que, na maioria das respostas, o estudante não situou o contexto de origem do conceito de indústria cultural.
Nas respostas corretas ( parcial ou integralmente ), de modo recorrente, foram encontradas as seguintes ideias :  a transformação dos bens culturais/artísticos em mercadoria industrializada, segundo algum interesse econômico relacionado à classe dominante ou às elites ;  a relação entre a modelagem da realidade pelos produtos culturais da indústria cultural, a alienação e a consciência massificada ;  a relação entre a modelagem da realidade pelos produtos culturais da indústria cultural e a manipulação das massas.
Houve, também, muitos casos de estudantes que abordaram a indústria cultural em diálogo com Walter Benjamin, a partir das ideias de perda da aura e da noção de reprodutibilidade técnica dos bens culturais.
Os principais erros e equívocos nas soluções dos itens ‘ a ’ e ‘ b ' estavam associados ao desconhecimento do tema proposto pelos frankfurtianos e do conceito de indústria cultural.
De modo recorrente, nestas respostas, os estudantes :  relacionaram a indústria cultural à produção da cultura popular ( culturais locais, do"povo"), o que beneficiaria os artistas populares e as economias regionais ;  definiram a indústria cultural como aquela responsável pela formação cultural das pessoas, segundo as prerrogativas de esclarecimento do iluminismo ;  definiram a indústria cultural como aquela responsável pelo enriquecimento dos saberes ;  definiram a indústria cultural como o tipo de indústria que produziria produtos feitos por máquinas e computadores e pelo avanço da tecnologia ;  definiram a indústria cultural como aquela responsável por gerar alegria, liberdade e esclarecimento.
As respostas consideradas nas análises a seguir são as que foram respondidas de modo correto e incorreto, segundo o padrão de resposta, e que foram pontuadas.
As respostas que foram desconsideradas, ou anuladas - conforme os critérios de correção - ou deixadas em branco, receberam nota zero, mas não entraram no cálculo dos percentuais aqui apresentados.
As respostas a que foram atribuídas notas baixas ( notas de zero a 30 pontos ) giraram em torno de 25 \% das respondidas pelos(as) estudantes do Bacharelado e em torno de 28 \% das respondidas pelos(as) estudantes da Licenciatura.
As tentativas de respostas que foram zeradas equivaleram a cerca de 9 \% em ambos os cursos.
Nas respostas dessa faixa, em geral, o/a estudante demonstrou desconhecer totalmente o assunto, e, por isso, ou confundiu conceitos, ou parafraseou o texto do enunciado na tentativa de escrever algo correto sobre o tema.
Foram também observados casos de textos com conteúdo aleatório e alguns escritos como um amontoado de palavras desconectadas.
Outras situações em que se justificaram as notas baixas foram : a fuga completa ao que foi solicitado pelo comando ; o texto desconexo, sem que o/a estudante conseguisse articular conceitualmente as duas partes do comando ; o texto como expressão do mais simplório senso comum sobre cultura, sem nenhuma análise filosófica que se espera de alguém que já integralizou, pelo menos, 70 \% dos créditos e da carga horária do curso.
As respostas consideradas medianas ( notas entre 40 e 70 pontos ) corresponderam a 34 \% das respondidas ( não se levando em conta as respostas em branco, as desconsideradas ou anuladas ), dos/das estudantes tanto do Bacharelado quanto da Licenciatura.
Em geral, o estudante demonstrou conhecer parcialmente o tema ao apresentar um texto com alguns argumentos presentes no padrão de resposta.
Foi mais comum haver respostas em que o/a estudante buscou, em resposta ao item ‘ a ’, definir a indústria cultural considerando a transformação dos bens culturais/artísticos em mercadoria industrializada, segundo algum interesse econômico relacionado à classe dominante ou às elites ; ou discutir a relação entre a modelagem da realidade pelos produtos culturais da indústria cultural, a alienação e a consciência massificada.
Em resposta ao item ‘ b ’, o mais recorrente foi apenas indicar ou desenvolver brevemente a relação entre a modelagem da realidade pelos produtos culturais da indústria cultural e a manipulação das massas.
Foi possível observar que, em parte significativa das respostas consideradas fracas e medianas, foram também identificadas dificuldades com o uso da norma culta da língua portuguesa, incluindo erros de coerência e coesão argumentativa e textual.
Tais fatos indicam a deficiência na apresentação do pensamento por escrito, em que pesem o cansaço decorrente do tempo grande de prova.
As respostas consideradas boas, com notas entre 80 e 100 pontos, equivaleram a cerca de 41 \% entre os/as estudantes do Bacharelado e a 38 \% entre os/as estudantes da Licenciatura.
Nelas, em geral, o estudante demonstrou conhecer o assunto e articulou, de forma completa ou parcial, os dois itens solicitados, conforme previsto no padrão de resposta.
Via de regra, esse tipo de resposta foi escrito de modo coerente e coeso, com demonstração de conhecimento seguro em relação ao que foi respondido.
Em boa parte das respostas consideradas boas, foram encontradas respostas criativas, quer pela extensão do texto – recheado de detalhes – quer pela coerência argumentativa e textual.
Tanto as respostas mais elaboradas quanto as menos elaboradas receberam boas notas por atenderem às expectativas assinaladas no padrão de resposta.
Pelos dados já apresentados, nota -se que não houve muita discrepância entre as notas das respostas dos formandos das duas habilitações, uma vez que estudantes de ambas não se diferenciaram quanto ao tipo de resposta.
Igualmente não foram percebidas diferenças significativas na construção da argumentação e na coerência e coesão textuais.
Nesse sentido, é preciso ressaltar que, em relação aos resultados dos exames do Enade anteriores, os/as estudantes da Licenciatura melhoraram consideravelmente o desempenho.
Foi frequente que respondentes da Licenciatura estabelecessem alguma relação com questões \textbf{pedagógicas} no corpo do texto, fato que demonstra preocupação em olhar o tema também pelo viés educacional.
Descartando -se o quantitativo de respostas anuladas ou desconsideradas ou em branco, é preciso ressaltar que 47 \% de estudantes do Bacharelado e 52 \% de estudantes da Licenciatura obtiveram notas entre zero e cinquenta ( 5,0 ), isto é, abaixo da média necessária para aprovação nas disciplinas dos cursos.
Houve um percentual alto de respostas em branco ( 15 \% no Bacharelado e 14 \% na Licenciatura ).
Nesse sentido, tendo em vista a importância e a disseminação do conceito de indústria cultural, e dada a possibilidade de que tenham tido acesso a esse conceito em diferentes disciplinas obrigatórias dos cursos de Filosofia, o desempenho dos/das estudantes na questão foi considerado mediano.
Notou -se que a elaboração dos textos das respostas foi, de modo geral, deficiente quanto à organização e à elaboração dos argumentos e quanto a erros gramaticais imperdoáveis em formandos e formandas de cursos de Filosofia, em que se pressupõem a leitura e a escrita como instrumentos consolidados nos últimos períodos.
Sugere -se que as coordenações dos cursos e o corpo docente como um todo devam considerar que a formação universitária precisa contribuir seriamente para desenvolver a capacidade de expressão escrita e de interpretação de textos filosóficos.
No entanto, é necessário registrar que, para que o aluno conseguisse responder no espaço destinado, era preciso ter capacidade de síntese e boa organização argumentativa.
Deve -se ressaltar, também, que a escrita manuscrita é feita por letras muitas vezes de difícil legibilidade, provavelmente como resultado do uso constante do teclado de um aparato técnico. 6.3.2.3 ANÁLISE DA QUESTÃO DISCURSIVA 4 DO COMPONENTE DE CONHECIMENTO ESPECÍFICO Questão anulada pela Comissão Assessora de Área. 6.3.2.4 COMENTÁRIOS SOBRE AS RESPOSTAS À QUESTÃO DISCURSIVA 4 Questão anulada pela Comissão Assessora da Área. 6.3.2.5 ANÁLISE DA QUESTÃO DISCURSIVA 5 DO COMPONENTE DE CONHECIMENTO ESPECÍFICO A Questão Discursiva 5 de Conhecimentos Específicos foi a mesma para as duas habilitações.
Os resultados para a Licenciatura constam da Tabela 6.19a e para o Bacharelado estão apresentados na Tabela 6.19b.
Cabe esclarecer que por terem respondido à mesma questão os resultados são comparáveis.
As Tabelas 6.19a e 6.19b contém, respectivamente para Licenciatura e Bacharelado, as informações relativas à Questão 5 do Componente de Conhecimento Específico.
Observa -se na Tabela 6.19a que o desempenho dos estudantes de todo o Brasil nesta questão foi muito inferior ao da Questão Discursiva 3.
A nota Média dos estudantes de todo o Brasil foi 17,4.
A maior Média foi registrada na região Sul ( 20,7 ), enquanto a menor Média foi registrada na região Nordeste ( 13,4 ).
Quanto à variabilidade das notas, o Desvio padrão dos estudantes do Brasil, como um todo, foi 24,3.
Enquanto o maior desvio foi encontrado na região Sul ( 26,3 ), o menor foi encontrado na região Norte ( 21,8 ).
A Mediana para o Brasil e em quase todas as regiões foi zero, com exceção da região Sul, onde a Mediana foi 10,0.
Para o conjunto de estudantes de Filosofia do Brasil, a nota Máxima foi 100,0, obtida por, pelo menos, um estudante em cada região, com exceção da região Norte que apresentou nota Máxima igual a 80,0.
A nota Mínima do Brasil foi zero, o mesmo valor encontrado em todas as regiões.
As Médias para os alunos de Licenciatura ( ver Tabela 6.19a ) foram, respectivamente para os de cursos Presenciais e a Distância, 19,6 e 15,0 ( 17,4 para o total ).
A nota Máxima ( 100,0 ) foi atingida por, pelo menos, um concluinte das regiões Nordeste, Sudeste, Sul e Centro-Oeste para a Modalidade Presencial, e das regiões Sudeste, Sul e Centro-Oeste para a Modalidade a Distância.
A Mediana foi zero para os cursos de ambas as Modalidades em todas as regiões, com exceção da Mediana dos cursos Presenciais da região Sudeste e Sul e dos cursos a Distância da região Centro-Oeste ( 25,0 ).
Cumpre notar que somente na região Norte não houve participação de estudantes de cursos a Distância de Licenciatura.
A nota Mínima foi zero para todas as combinações de Modalidade de Ensino e região.
Na Tabela 6.19b, constam as informações relativas à Questão 5 do conjunto do Componente de Conhecimento Específico.
O desempenho dos estudantes de todo o Brasil nesta questão foi muito inferior ao verificado na Questão 3.
A nota Média dos estudantes de todo o Brasil foi 18,6.
A maior Média foi registrada na região Sul ( 21,6 ), enquanto a menor Média foi registrada na região Norte ( 11,1 ).
Quanto à variabilidade das notas, o Desvio padrão dos estudantes do Brasil, como um todo, foi 25,1.
Enquanto o maior desvio foi encontrado na região Sul ( 26,2 ), o menor foi encontrado na região Centro-Oeste ( 21,9 ).
A Mediana para o Brasil foi 10,0, assim como as Medianas de duas regiões.
Nas regiões Norte, Nordeste e Centro-Oeste as Medianas foram 0,0.
Para o conjunto de estudantes de Filosofia do Brasil, a nota Máxima foi 100,0, obtida por, pelo menos, um estudante de todas as regiões.
A nota Mínima do Brasil foi zero, o mesmo valor encontrado em todas as regiões.
As Médias para os alunos de Bacharelado ( ver Tabela 6.19b ) foram, respectivamente para os de cursos Presenciais e a Distância, 18,8 e 17,4 ( 18,6 para o total ).
A nota Máxima ( 100,0 ) foi atingida por, pelo menos, um concluinte em todas as regiões e em ambas Modalidades de Ensino, com exceção dos concluintes da região Centro-Oeste na Modalidade Presencial, onde a Máxima foi 70,0.
A Mediana foi zero para os cursos de ambas as Modalidades em todas as regiões, com exceção da Mediana do curso Presencial das regiões Sudeste e Sul ( 10,0 ).
Cumpre notar que somente nas regiões Norte e Nordeste não houveram participação de estudantes de cursos a Distância de Bacharelado.
A nota Mínima foi zero para todas as combinações de Modalidade de Ensino e região.
No Gráfico 6.12, é representada a distribuição das notas da Questão Discursiva 5 do Componente de Conhecimento Específico.
Para os alunos de Licenciatura, essa distribuição tem moda na categoria de notas zero, daqueles que apresentaram resposta à questão ( 24,3 \% ).
Chega a 52,3 \% a soma desse contingente com o daqueles que deixaram a resposta à questão em branco.
Excetuando as notas zero, o intervalo ( 0 ; 10 ], que corresponde a 12,4 \% dos respondentes, é o de maior frequência.
Dentre os alunos de Bacharelado, é representada a distribuição das notas da Questão Discursiva 5 do Componente de Conhecimento Específico.
Essa distribuição tem moda no intervalo ( 0 ; 10 ], correspondendo a 11,8 \% dos respondentes.
O percentual dos participantes que tiraram nota zero, tendo respondido a essa questão, ou por terem deixado a resposta em branco chega a 48,9 \%. 6.3.2.6 COMENTÁRIOS SOBRE AS RESPOSTAS À QUESTÃO DISCURSIVA 5 A Questão Discursiva 5 foi proposta para estudantes do Bacharelado e da Licenciatura em Filosofia, em que foi abordada a contribuição de Karl Popper para a Filosofia da Ciência.
Nas Diretrizes Curriculares Nacionais para os cursos de Filosofia estão destacadas as disciplinas Filosofia da Ciência ou Epistemologia e Teoria do Conhecimento como relevantes para a estruturação do conhecimento filosófico a ser desenvolvido durante a graduação.
O perfil do egresso indica a necessidade de, nas graduações, ser desenvolvida a capacidade dos/das estudantes quanto à contextualização das ideias filosóficas, à proficiência na argumentação e na exposição dessas ideias, à compreensão das ciências e dos saberes no contexto das culturas, e quanto à capacidade de interpretarem e comentarem filosoficamente a produção científica e a cientificidade.
Nesse sentido, a questão se encaixa, a princípio, no escopo das graduações em Filosofia, nas quais aparecem as discussões sobre o caráter da cientificidade e a demarcação do que vem a ser o conhecimento científico, sobre os métodos de investigação, sobre a dedução e a indução, sendo a perspectiva de Karl Popper uma das abordagens possíveis para se considerarem esses problemas.
Na questão, eram apresentados dois textos.
No primeiro, constituído por um pequeno trecho da obra"A Lógica da Pesquisa Científica", de Popper, era feita uma referência famosa aos cisnes brancos.
No segundo texto, havia um comentário sobre o critério da falseabilidade.
Na questão, solicitava -se que o respondente explicasse como a abordagem de Popper demarcava uma teoria científica e o que não é científico e que relacionasse essa demarcação ao problema da indução.
A questão era clara e exigia que a explicação dada pelo estudante fosse feita segundo a abordagem popperiana.
O padrão de resposta contemplava boa diversidade de argumentos que poderiam ser utilizados pelos estudantes.
Para efeito de pontuação, a questão foi subdividida em dois itens, a fim de valorar adequadamente a diversidade de caminhos argumentativos utilizados pelos respondentes.
As respostas foram pontuadas levando -se em consideração, principalmente, a abordagem da perspectiva de K.
Popper acerca da distinção entre teorias científicas e não científicas, a partir do critério popperiano da falseabilidade ou pela crítica ao verificacionismo ou pela relação com qualquer método científico ou, ainda, pela relação de uma teoria científica com a provisoriedade/historicidade.
Quanto ao problema da indução, esperava -se que o estudante desenvolvesse sua argumentação a partir da relação entre indução e verificação ou da relação entre cientificidade e verificação.
Apesar das orientações das Diretrizes, as disciplinas “ Filosofia da Ciência ” ou “ Epistemologia ” não possuem caráter obrigatório em todos os cursos de Filosofia, o que tornou a questão difícil para boa parte dos formandos.
Os/as estudantes que conseguiram responder corretamente ( parcial ou integralmente ) argumentaram usando uma das possibilidades presentes no padrão de resposta, que contemplou boa diversidade de caminhos possíveis.
As respostas incorretas resultaram do desconhecimento do assunto, segundo a abordagem de K.
Popper, ou da compreensão do que demarca uma teoria científica, ou ainda da confusão entre dedução e indução, ou da confusão entre falseabilidade e falsificação de resultados.
A qualidade e a adequação da linguagem adotada pelos estudantes na redação da resposta variou bastante, sem associação com o tipo de habilitação : Bacharelado ou Licenciatura.
O percentual de notas entre zero e cinquenta foi alto, em torno de 86 \% no Bacharelado e de 88 \% na Licenciatura.
Dentre os estudantes de Bacharelado e de Licenciatura que receberam nota zero na Questão 5, respectivamente, 30 \% e 35 \% estavam em branco ou foram desconsideradas ou nulas.
Parte considerável dos/das estudantes com respostas nessa faixa de notas redigiram textos em linguagem bem primária e simplória, próxima do senso comum, demonstraram dificuldade na construção de argumentos, não tendo conseguido expor seus conhecimentos com alguma proficiência.
Na correção, ficou evidente que a maioria dos/das estudantes enfrentou dificuldades para responder à questão, corroborando a percepção prévia de que seria uma questão difícil.
Em geral, os/as estudantes que deram respostas corretas parcial ou integralmente lançaram mão de argumentos presentes no padrão de resposta quanto à : a ) explicação do critério de falseabillidade e sua relação com a cientificidade e com o problema da indução ; b ) crítica de Popper ao verificacionismo ; c ) provisoriedade/historicidade das verdades científicas ; d ) caracterização de uma teoria científica pela presença de algum método científico ; e ) explicação do problema presente na indução.
Os principais erros e equívocos percebidos nas soluções apresentadas pelos estudantes decorreram de : a ) confusão entre dedução e indução, fato que não se justifica em cursos de Filosofia, pois desde a análise da Filosofia na Grécia antiga, esses dois temas estão presentes ; b ) incompreensão dos limites da indução, fato que também não se justifica em cursos de Filosofia, pois, mesmo em conteúdos da disciplina Filosofia Moderna, há referências às dificuldades ligadas à indução ; c ) incompreensão ou, pelo menos, falta de posicionamento fundamentado em relação ao que demarca as verdades científicas, fato que igualmente não se justifica em cursos de Filosofia, uma vez que não se estuda pensamento moderno sem lidar com o surgimento da ciência moderna e com a questão do método ; d ) confusão entre falseabilidade e falsificação de resultados, fato que demonstra total desconhecimento da abordagem popperiana.
Foi comum os/as estudantes usarem o recurso da paráfrase dos dois textos do enunciado para a construção de argumentos ou, ainda, apenas indicar um caminho de resposta sem desenvolvimento de uma argumentação.
As respostas que obtiveram notas baixas ( zero até 30 pontos ) constituíram 49 \% das corrigidas dos/das estudantes do Bacharelado e 44 \% das corrigidas dos/das estudantes da Licenciatura, não levando em conta as respostas em branco ou que foram anuladas ou desconsideradas.
As respostas corrigidas e zeradas equivaleram a cerca de 20 \% em ambos os cursos.
Nessas respostas, em geral, o/a estudante demonstrou desconhecer o assunto, e, por isso, ou confundiu conceitos, ou apenas parafraseou o texto do enunciado.
Outras situações em que se justificaram as notas baixas foram : a fuga completa ao solicitado no comando ; texto desconexo sem que o/a estudante conseguisse articular conceitualmente o que foi solicitado na questão ; texto como expressão do mais simplório senso comum sobre ciência, verdade científica, falseabilidade, sem nenhuma análise filosófica que se espera de alguém que já integralizou pelo menos 70 \% dos créditos e da carga horária do curso.
As respostas consideradas medianas ( entre 40 e 70 pontos ), em ambas as habilitações, corresponderam a, aproximadamente, 20 \% das corrigidas ( não se levando em conta as respostas em branco, as desconsideradas ou anuladas ).
Nas respostas dessa faixa, em geral, o estudante demonstrou conhecer parcialmente o que foi solicitado ao abordar o tema a partir de algum dos argumentos presentes no padrão de resposta.
As respostas consideradas boas, com notas entre 80 e 100 pontos, equivaleram a 1,7 \% do total de participantes do Bacharelado e a 1,3 \% dos da Licenciatura.
Nelas, em geral, o estudante demonstrou conhecer o assunto e articulou, de forma completa ou parcial, os dois itens pedidos, conforme previsto pelo padrão de resposta.
Pelo resultado, ficou evidente o desconhecimento do assunto entre parte considerável dos/as discentes do Bacharelado e da Licenciatura em Filosofia.
Pelos dados expostos, percebe -se que não existiu grande diferença entre as notas dos formandos das duas habilitações.
Igualmente, não foram percebidas diferenças significativas na construção da argumentação e na capacidade de redigir um texto com coerência e coesão.
Cabe observar que nas respostas de estudantes da Licenciatura foi comum eles estabelecerem alguma relação com questões \textbf{pedagógicas} no corpo do texto, fato que demonstra preocupação em olhar o tema também pelo viés educacional.
Em parte considerável das respostas consideradas fracas e medianas, os estudantes, além de demonstrarem problemas conceituais, cometeram erros de coerência e coesão argumentativa e textual ( conhecimento da norma culta da língua portuguesa ), o que indica deficiência na apresentação do pensamento por escrito.
Notou -se que a elaboração dos textos das respostas foi, de modo geral, deficiente quanto à organização e à construção dos argumentos e quanto a erros gramaticais que não deveriam fazer parte de textos redigidos por formandos e formandas de cursos de Filosofia, já que, nos cursos, pressupõe -se que haja grande volume de leitura e de escrita durante a formação universitária.
Sugere -se que os cursos e o corpo docente como um todo possam contribuir mais para o desenvolvimento da capacidade de expressão escrita e de interpretação de textos filosóficos ou não.
Na tentativa de se entenderem os possíveis motivos para o fraco desempenho dos/das estudantes na Questão 5, foi realizado um levantamento dos projetos \textbf{pedagógicos} de importantes cursos de Filosofia ( em torno de 30 \% dos cursos do país ).
Apesar da orientação das Diretrizes Nacionais para os cursos de Filosofia em relação a temas, problemas e abordagens relativos às produções científicas, tal como exposto no início desta análise, foi constatado que a disciplina Filosofia da Ciência ( ou Epistemologia ) não é obrigatória em todos os cursos.
Na maioria dos cursos, a disciplina é optativa ou eletiva condicionada.
Ao se analisarem os programas da disciplina Filosofia da Ciência ou Epistemologia ( obrigatória ou optativa/eletiva ) foram constatadas situações em que as obras de Karl Popper não constam da bibliografia básica, ou da bibliografia complementar ou mesmo do programa da disciplina.
Provavelmente, o conjunto desses aspectos foi decisivo para o fraco desempenho dos/das estudantes nessa questão discursiva.
Devido à relevância da reflexão crítica acerca do caráter de cientificidade e das verdades científicas, especialmente em tempos de negacionismos, sugere -se maior cuidado com os temas, problemas e abordagens relacionados às produções científicas e às ciências.
Como na Questão 3, para que os alunos conseguissem responder no espaço de 15 linhas, era preciso ter capacidade de síntese e boa organização argumentativa, especialmente em textos manuscritos.
Vale lembrar, ainda, que no texto manuscrito, frequentemente, se encontram problemas de caligrafia, o que pode ser resultado do uso constante do teclado de um aparato tecnológico. 6.3.2.7 CONSIDERAÇÕES FINAIS SOBRE AS QUESTÕES DISCURSIVAS DO COMPONENTE DE CONHECIMENTO ESPECÍFICO Nesse relatório, foram apresentadas as considerações sobre as respostas dos/das estudantes do Bacharelado e da Licenciatura em Filosofia, nas questões discursivas de número 3 e de número 5 de Conhecimentos Específicos.
A Questão de número 4 foi anulada pela Comissão Assessora da Área, para as duas habilitações.
O desempenho dos/das estudantes na Questão Discursiva 3, em que se abordou um tema que costuma ser tratado em diferentes disciplinas do curso de Filosofia, foi bem superior ao observado na Questão 5.
Supõe -se que o desempenho nesta questão tenha sido prejudicado pelo fato de a disciplina Filosofia da Ciência não ser obrigatória em todos os cursos de Bacharelado e de Licenciatura em Filosofia, tal como explicado na análise da questão.
Foi possível observar, em ambas as questões específicas, que, nas respostas consideradas fracas e medianas, foram também percebidos erros de coerência e coesão argumentativa e textual ( conhecimento da norma culta da língua portuguesa ).
Em linhas gerais, notou -se que nos textos redigidos pelos estudantes, houve ocorrência de deficiências quanto à organização e à elaboração de argumentos, além da presença de erros gramaticais que não cabem a formandos e formandas de cursos de Filosofia.
Essa dificuldade põe em xeque a formação para a realização das competências e habilidades previstas pelas Diretrizes Nacionais da Área para os egressos do Bacharelado e da Licenciatura em Filosofia.
Sugere -se que as coordenações dos cursos e o corpo docente como um todo planejem formas de contribuir mais eficientemente para o desenvolvimento das capacidades de expressão escrita, de construção de argumentos e de interpretação de textos.

% matched lemmas: crescer, pedagógico, recorte
\end{textsample}
